%======================================================================
\section{Area-preserving disc maps from Morse charts}\label{sec:pln}

The fibre map of the handle construction (\Msec{sec:hdl}) must be the endpoint of a compactly
supported Hamiltonian isotopy of a disc which is the identity near the boundary and satisfies a comparison
inequality. In \cite{OpenAI2026} such an isotopy is obtained from a rearrangement lemma by nested discs, a construction it
credits to \cite[Lemma~3.2]{SchlenkFolding}, and Moser's method. This section constructs it explicitly. The construction is used twice: for the explicit
inner map for equal balls (\S\ref{sec:pll}) and for the realisation of coordinate rearrangements of convex
functions (\S\ref{sec:plc}). The input is a pair of potentials, $U$ in a moment variable $p$ and $V$ in an angle $\theta$, each with a single
nondegenerate minimum near which it is exactly quadratic; in the slice family both depend smoothly on a parameter $y$
(the other moments).
The output is a compactly supported Hamiltonian isotopy of the plane which preserves a disc of prescribed
area, is the identity on an outer collar of it, and whose endpoint on a prescribed core disc is an explicit
area-preserving map carrying round circles onto level sets of $U+V$. The two families of potentials are defined later, in \S\ref{sec:pll} and in
Definition~\ref{def:plc-data}; this section uses only the properties listed in Definition~\ref{def:pln-potential-pair}, and a
reader of the equal-ball case can ignore every statement about the slice family.

The construction has two halves. The first (\S\S\ref{sec:pln-morse}--\ref{sec:pln-inner}) produces the
endpoint. Signed square roots of the potentials give a Morse chart. A rotation normalisation on each circle
and a radial area normalisation make it area preserving, and the polar chart places it in the disc
(Theorem~\ref{thm:pln-inner}). The second half (\S\S\ref{sec:pln-charts}--\ref{sec:pln-isotopy}) joins the
identity to this endpoint without changing it on the core. Interpolating the Morse chart linearly with its
linearisation, and rescaling the polar chart, gives a path of symplectic embeddings. A five-segment path of
embeddings, all with image in a fixed smaller disc, connects it to the identity. The ambient flow of this
path, cut off in the source variable, is made area preserving by one explicit normalisation of the whole
path. The Hamiltonian is the radial primitive of the resulting flow, normalised at a fixed point of the
collar (Theorems~\ref{thm:pln-isotopy} and~\ref{thm:pln-param}).

\subsection*{Conventions}
For $P>0$, $D(P)=\{z\in\C:\pi|z|^2<P\}$ is the open disc of \emph{area} $P$ and $\bar D(P)$ its
closure; $D(\infty)=\C$. Every plane carries its standard area form: $\omega_\C$ on the source and target
discs $\C$, $du\wedge dv$ on the Morse plane $\R^2\ni w=(u,v)=u+iv$, and $dx\wedge d\theta$ on the
rectangle coordinates $(x,\theta)$. A map between open subsets of planes is \emph{symplectic} if its
Jacobian determinant is identically $1$. In polar area coordinates $p=\pi|z|^2$,
$\vartheta=\arg(z)/2\pi\in\R/\Z$ (the fibre angle $\theta_j$ of the main text) one has
$\omega_\C=dp\wedge d\vartheta$ on $\C\setminus\{0\}$; under the polar chart $C_y$ of \S\ref{sec:pln-inner}
the rectangle angle is $\theta=\vartheta-\frac12$. $\mathrm{Rot}_\alpha$ is the rotation by the angle $\alpha$
(radians). $S$ is the flat step of Definition~\ref{def:pre-flat}; by Lemma~\ref{lem:pre-flat} it is
$C^\infty$ and nondecreasing, $S=0$ on $(-\infty,0]$, $S=1$ on $[1,\infty)$, and all derivatives of $S$
vanish at $0$ and $1$.

In this section $E$ always denotes an \emph{energy level} (a value of a potential), $J$ a positive
density on a plane, and $h>0$ the area of the disc $D(h)$ preserved by the isotopy. A function on an
arbitrary subset $\Omega\subset\R^N$ is called \emph{smooth} if every point of $\Omega$ has an open
neighbourhood in $\R^N$ on which it is the restriction of a $C^\infty$ function; on open sets this is the
usual notion, and at the endpoints of a parameter interval it means smoothness with one-sided derivatives.
Fibre Hamiltonians follow the convention $\iota_X\omega_\C=-dQ$ of \Msec{sec:cmp}.

%----------------------------------------------------------------------
\subsection{Admissible potential pairs}\label{sec:pln-pairs}

\begin{definition}[Admissible potential pair]\label{def:pln-potential-pair}
An \emph{admissible potential pair} is one of the following two families: the modified equal-ball wells of
\S\ref{sec:pll} ($k\ge10$, $\frac78<t_0<1$; Lemma~\ref{lem:pll-pair}), or the modified slice wells of
Definition~\ref{def:plc-data} over the open parameter set $\mathcal O$ of the planar data (Lemma~\ref{lem:plc-data}).
For either family let $\mathcal Y\subset\R^d$ be its open parameter domain ($d=0$ meaning a single point), $h>0$ its
disc area, and $U_y(p)$, $V_y(\theta)$ its potentials on open intervals $I_U,I_V$. The following properties are proved
for both families in the lemmas just cited:
\begin{enumerate}[label=(P\arabic*)]
\item\label{it:pln-P1} The potentials are jointly smooth. There is a unique
$p^*(y)\in I_U\cap(0,h)$ with $U_y'(p)(p-p^*(y))>0$ away from the minimum,
$U_y''(p^*(y))>0$, and $V_y(0)=0$, $V_y'(\theta)\theta>0$ away from zero,
$V_y''(0)>0$.
\item\label{it:pln-P2} Both potentials tend to $+\infty$ at their interval
endpoints. Their centred wells are exactly quadratic on neighbourhoods of zero;
the widths are positive uniformly on the parameter domain used in a construction.
The quadratic coefficients are smooth positive functions of the parameters.
\end{enumerate}
We write $U_{\min}(y)=U_y(p^*(y))$. The pair is \emph{exactly quadratic near the minimum} if there are
$\eps_U,\eps_V>0$ with $U_y(p)-U_{\min}(y)=\frac12U_y''(p^*(y))(p-p^*(y))^2$ for $|p-p^*(y)|\le\eps_U$ and
$V_y(\theta)=\frac12V_y''(0)\theta^2$ for $|\theta|\le\eps_V$.
\end{definition}

The modified wells need not be convex; only the signs of their derivatives and their quadratic germs are used.

%----------------------------------------------------------------------
\subsection{Morse inverse charts}\label{sec:pln-morse}

\begin{lemma}[Morse inverse charts]\label{lem:pln-morse-inverse}
\lean{SymplecticPacking.Inner.SigmaR_smooth, SymplecticPacking.Inner.Xtilde_smooth,
 SymplecticPacking.Inner.Rtilde_Xtilde, SymplecticPacking.Inner.Xtilde_linear,
 SymplecticPacking.Inner.Ytilde_smooth, SymplecticPacking.Inner.Vtilde_Ytilde,
 SymplecticPacking.Inner.Ytilde_linear, SymplecticPacking.Branch2.Planar.MorsePair.Fmor_det,
 SymplecticPacking.Inner.Fmor_level, SymplecticPacking.Branch2.Planar.MorsePair.Acal_eq_area}
\uses{def:pln-potential-pair}
Let $(U,V)$ be an admissible potential pair over $\mathcal Y$.
\begin{enumerate}[label=(\roman*)]
\item The minimum point $y\mapsto p^*(y)$ and $y\mapsto U_{\min}(y)$ are $C^\infty$ on $\mathcal Y$.
\item On the open set $\mathcal D_U=\{(y,x):y\in\mathcal Y,\ p^*(y)+x\in I_U\}$ put
 \[
  \Sigma_U(y,x)=\operatorname{sign}(x)\sqrt{U_y(p^*(y)+x)-U_{\min}(y)}.
 \]
 This is the signed square root, with its removable singularity at $x=0$; and $\Sigma_U$ is $C^\infty$ with
 $\partial_x\Sigma_U>0$. For each $y$, $\Sigma_U(y,\cdot)$ is an increasing diffeomorphism of
 $I_U-p^*(y)$ onto $\R$; its inverse $X_y$ is $C^\infty$ jointly in $(y,u)\in\mathcal Y\times\R$, and
 \begin{gather*}
  U_y\bigl(p^*(y)+X_y(u)\bigr)-U_{\min}(y)=u^2,\\
  \operatorname{sign}X_y(u)=\operatorname{sign}u,\qquad
  X_y'(0)=a_0(y):=\sqrt{2/U_y''(p^*(y))} .
 \end{gather*}
\item In the same way $\Sigma_V(y,\theta)=\operatorname{sign}(\theta)\sqrt{V_y(\theta)}$
 is an increasing diffeomorphism of $I_V$ onto $\R$ for each $y$, and its inverse $Y_y$ is $C^\infty$ in
 $(y,v)$, with $V_y(Y_y(v))=v^2$, $\operatorname{sign}Y_y(v)=\operatorname{sign}v$ and
 $Y_y'(0)=b_0(y):=\sqrt{2/V_y''(0)}$.
\item The Morse chart $F_y(u,v)=(X_y(u),Y_y(v))$ is a $C^\infty$ diffeomorphism of $\R^2$ onto the
 \emph{Morse rectangle} $\mathcal R_y=(I_U-p^*(y))\times I_V$, jointly smooth in $(y,w)$, with
 \[
  F_y^*(dx\wedge d\theta)=J_y\,du\wedge dv,\qquad J_y(u,v)=X_y'(u)\,Y_y'(v)>0,\qquad J_y(0)=j_0(y):=a_0(y)b_0(y),
 \]
 and $U_y(p^*(y)+x)-U_{\min}(y)+V_y(\theta)=u^2+v^2$ for $(x,\theta)=F_y(u,v)$. Consequently, for $E\ge0$,
 \[
  \int_{|w|^2\le E}J_y(w)\,dw=\operatorname{area}\bigl\{(p,\theta)\in I_U\times I_V:\ U_y(p)-U_{\min}(y)+V_y(\theta)\le E\bigr\},
 \]
 and these areas increase to $|I_U|\,|I_V|\in(0,\infty]$ as $E\to\infty$.
\item If the pair is exactly quadratic near the minimum, then $X_y(u)=a_0(y)u$ for $|u|\le\eps_U/a_0(y)$
 and $Y_y(v)=b_0(y)v$ for $|v|\le\eps_V/b_0(y)$; hence $J_y=j_0(y)$ on the disc $|w|\le\rho_T(y)$,
 $\rho_T(y)=\min(\eps_U/a_0(y),\eps_V/b_0(y))$.
\end{enumerate}
\end{lemma}

\begin{proof}
(i) $\partial_pU_y'(p)=U_y''(p)>0$ at $p=p^*(y)$, so the implicit function theorem gives, near each
$y_0$, a smooth solution of $U_y'(p)=0$ through $p^*(y_0)$; by \ref{it:pln-P1} the zero of $U_y'$ is
unique, so this solution is $p^*(y)$. Then $U_{\min}(y)=U_y(p^*(y))$ is smooth.

(ii) $\mathcal D_U$ is open since $p^*$ is continuous. The quadratic
germ gives $\Sigma_U(y,x)=\sqrt{U_y''(p^*)/2}\,x$ on a common neighbourhood
of $x=0$, which is jointly smooth there. Away from $x=0$, positivity of
$U_y(p^*+x)-U_{\min}$ gives smoothness directly from the signed square-root
formula. For $x\ne0$,
differentiating $\Sigma_U^2=U_y(p^*+x)-U_{\min}$ gives $2\Sigma_U\partial_x\Sigma_U=U_y'(p^*+x)$; both
$\Sigma_U$ and $U_y'(p^*+x)$ have the sign of $x$, so $\partial_x\Sigma_U>0$; at $x=0$,
$\partial_x\Sigma_U=\sqrt{U_y''(p^*)/2}>0$. By \ref{it:pln-P2}, $\Sigma_U(y,\cdot)$
tends to $\mp\infty$ at the two endpoints of $I_U-p^*(y)$, so it is an increasing bijection onto $\R$.
The map $(y,x)\mapsto(y,\Sigma_U(y,x))$ is a smooth bijection of $\mathcal D_U$ onto
$\mathcal Y\times\R$ with invertible (triangular) derivative; by the inverse function theorem its local
smooth inverses exist everywhere, and by injectivity they agree with the global inverse
$(y,u)\mapsto(y,X_y(u))$, which is therefore smooth. The identities follow from
$\Sigma_U(y,X_y(u))=u$, and $X_y'(0)=1/\partial_x\Sigma_U(y,0)=\sqrt{2/U_y''(p^*)}$.

(iii) is (ii) for $V$, whose minimum point is $0$ and minimum value $0$.

(iv) $F_y$ is a product of diffeomorphisms; its Jacobian is $X_y'Y_y'>0$, and $J_y(0)=a_0b_0$. The level
identity is the sum of the identities in (ii) and (iii). The change of variables formula applied to
$F_y$ on $\{|w|^2\le E\}$ gives the area formula, after the translation $p=p^*(y)+x$. As $E\to\infty$ the
discs exhaust $\R^2$, so their images exhaust $\mathcal R_y$ and the areas increase to
$\operatorname{area}\mathcal R_y=|I_U||I_V|$.

(v) The exact quadratic formula on $|x|\le\eps_U$ gives $\Sigma_U(y,x)=x/a_0$; hence $X_y(u)=a_0u$ for $|u|\le\eps_U/a_0$.
The same holds for $V$. For $|w|\le\rho_T$ both $|u|$ and $|v|$ are at most $\rho_T$, so
$J_y(w)=a_0b_0$.
\end{proof}

%----------------------------------------------------------------------
\subsection{Rotation and area normalisation}\label{sec:pln-normalisation-maps}

The two normalisations are defined for the positive densities used in the construction. Throughout this
subsection we fix the following setting, which we call \emph{setting (N)}: $\mathcal P\subset\R^N$ is
an arbitrary subset (the parameters), $\mathcal U\subset\R^N\times\R^2$ is an open set containing
$\mathcal P\times\R^2$, and $J\colon\mathcal U\to\R$, $(\sigma,w)\mapsto J_\sigma(w)$, is $C^\infty$ and
positive on $\mathcal P\times\R^2$, and $J$ has a constant germ at $w=0$: locally in the parameter there are
$\rho_0>0$ and a smooth positive function $j(\sigma)$ with $J_\sigma(w)=j(\sigma)$ for $|w|\le\rho_0$. The densities
used are the Morse densities of the two potential families and their interpolations (\S\ref{sec:pln-charts}), the
Morse density of \S\ref{sec:scalar-cmp}, and the whole-path densities of \S\ref{sec:pln-normal}. Points of $\R^2$ are
identified with complex numbers, so that $w\zeta$ is the complex product.

\begin{definition}[Rotation and area normalisation]\label{def:pln-normalisations}
\lean{SymplecticPacking.Inner.Jbar, SymplecticPacking.Inner.Bfun, SymplecticPacking.Inner.Trot,
 SymplecticPacking.Inner.Acal, SymplecticPacking.Inner.Einv, SymplecticPacking.Inner.Darea}
In setting (N), for $\sigma\in\mathcal P$, $w\in\R^2$ and $E\ge0$ put
\begin{gather*}
 \bar J_\sigma(|w|^2)=\frac1{2\pi}\int_0^{2\pi}J_\sigma(\mathrm{Rot}_\alpha w)\,d\alpha,\qquad
 \mathcal M_\sigma(w)=\frac1{2\pi}\int_0^{2\pi}(\alpha-\pi)\,J_\sigma(\mathrm{Rot}_\alpha w)\,d\alpha,\\
 T_\sigma(w)=\mathrm{Rot}_{\tau_\sigma(w)}\,w,\qquad \tau_\sigma(w)=\mathcal M_\sigma(w)/\bar J_\sigma(|w|^2),\\
 \mathcal A_\sigma(E)=\int_{|w|^2\le E}J_\sigma(w)\,dw,\qquad
 \mathcal A_\sigma(\infty)=\lim_{E\to\infty}\mathcal A_\sigma(E)\in(0,\infty],\\
 \mathcal G_\sigma(w)=\frac1\pi\int_{|\zeta|\le1}J_\sigma(w\zeta)\,d\zeta .
\end{gather*}
Let $E_\sigma\colon[0,\mathcal A_\sigma(\infty))\to[0,\infty)$ be the inverse of $\mathcal A_\sigma$
(Lemma~\ref{lem:pln-area}) and, for $z\in D(\mathcal A_\sigma(\infty))$ with $I=\pi|z|^2$,
\[
 D_\sigma(z)=\sqrt{\pi E_\sigma(I)/I}\;z\quad(z\ne0),\qquad D_\sigma(0)=0 .
\]
$T_\sigma$ is the \emph{rotation normalisation} and $D_\sigma$ the \emph{area normalisation} of $J_\sigma$.
When a single density $J$ is meant we write $\bar J$, $\mathcal M_J$, $T_J$, $\mathcal A_J$, $E_J$, $\mathcal G_J$, $D_J$.
\end{definition}

The normalisations of $J_\sigma$ at a point $w$ with $|w|\le R$, or at a point $z$ with $E_\sigma(\pi|z|^2)\le R^2$,
depend only on $J_\sigma$ on the disc $|w|\le R$: the circle and disc integrals defining $\bar J_\sigma$,
$\mathcal M_\sigma$ and $\mathcal A_\sigma$ on $[0,R^2]$ stay in that disc. Hence on such points they coincide with the
normalisations of
\[
 J^{\rm ext}_\sigma(w)=\vartheta_R(w)J_\sigma(w)
                  +(1-\vartheta_R(w))j(\sigma),\qquad
 \vartheta_R(w)=1-S\left(\frac{|w|^2-R^2}{3R^2}\right),
\]
which is smooth, positive, has the same germ, equals $J_\sigma$ for $|w|\le R$, and has
$\mathcal A_{J^{\rm ext}_\sigma}(\infty)=\infty$. For the slice family we use this with $R^2=\bar E$
(Definition~\ref{def:plc-data}), so that the source disc $D(p^\dagger)$ is covered by \ref{it:pln-D5}; it is a device
of proof only, and all maps are defined with the original densities.

On the circle of radius $\rho$ the density defines a measure of total mass $2\pi\rho\bar J_\sigma(\rho^2)$.
On each such circle $T_\sigma$ is determined up to a rotation by the condition that it carries this measure to
the uniform measure of the same mass, i.e.\ by the derivative of its angle (Lemma~\ref{lem:pln-rotation}(ii)). The
rotation is fixed by the weighted mean $\mathcal M_\sigma$; this choice needs no base point on the circle and is a
Cartesian function of $w$. Likewise $\mathcal G_\sigma$ writes the
mean of $J_\sigma$ over the disc of radius $|w|$ as an integral over a fixed disc, which is what makes the
area normalisation smooth at the origin.

\begin{lemma}[Rotation normalisation]\label{lem:pln-rotation}
\lean{SymplecticPacking.Inner.Bfun_deriv_arg, SymplecticPacking.Inner.norm_Trot,
 SymplecticPacking.Inner.Trot_bijective, SymplecticPacking.Inner.Trot_det,
 SymplecticPacking.Inner.Trot_eq_id_of_const, SymplecticPacking.Inner.Trot_param_smooth,
 SymplecticPacking.Inner.TrotInv_param_smooth}
\uses{def:pln-normalisations}
In setting (N), let $\sigma\in\mathcal P$.
\begin{enumerate}[label=(\roman*)]
\item $\partial_{\arg}\mathcal M_\sigma=J_\sigma-\bar J_\sigma(|w|^2)$, where $\partial_{\arg}$ is the
 derivative along $\psi$ in $w=\rho e^{i\psi}$.
\item $T_\sigma$ maps each circle $\{|w|=\rho\}$ onto itself; for $\rho>0$, writing
 $T_\sigma(\rho e^{i\psi})=\rho e^{i\psi^\sharp(\psi)}$, one has $(\psi^\sharp)'=J_\sigma/\bar J_\sigma>0$ and
 $\psi^\sharp(\psi+2\pi)=\psi^\sharp(\psi)+2\pi$. $T_\sigma$ is a bijection of $\R^2$ and $T_\sigma(0)=0$.
\item The maps $(\sigma,w)\mapsto T_\sigma(w)$ and $(\sigma,w)\mapsto T_\sigma^{-1}(w)$ are smooth on
 $\mathcal P\times\R^2$, each $T_\sigma$ is a diffeomorphism of $\R^2$ with $DT_\sigma(0)=\mathrm{Id}$, and
 $T_\sigma^*\bigl(\bar J_\sigma(|w|^2)\,du\wedge dv\bigr)=J_\sigma\,du\wedge dv$ on all of $\R^2$.
\item If $J_\sigma(w)=j$ for $|w|\le\rho_0$, then $\mathcal M_\sigma=0$ and $T_\sigma=T_\sigma^{-1}=\mathrm{id}$
 on $\{|w|\le\rho_0\}$.
\end{enumerate}
\end{lemma}

\begin{proof}
\emph{Local formulas.} Let $\mathcal U'=\mathcal U\cap\{J>0\}$, an open set containing
$\mathcal P\times\R^2$. Given $\sigma_0\in\mathcal P$ and $\rho_1>0$, the compact set
$\{\sigma_0\}\times\bar B(0,\rho_1)$ lies in $\mathcal U'$, so there is an open ball $N_0\ni\sigma_0$ with
$N_0\times B(0,\rho_1)\subset\mathcal U'$. On $N_0\times B(0,\rho_1)$, which is invariant under the
rotations of the second factor, the integrands defining $\bar J_\sigma(|w|^2)$ and $\mathcal M_\sigma(w)$ are
$C^\infty$ in $(\alpha,\sigma,w)\in[0,2\pi]\times N_0\times B(0,\rho_1)$ and positive, respectively smooth, so
differentiation under the integral sign shows that both formulas define $C^\infty$ functions there, with
$\bar J>0$. Hence the formula for $T_\sigma$ defines a $C^\infty$ map $\tilde T$ on
$N_0\times B(0,\rho_1)$, which is a local smooth extension of $(\sigma,w)\mapsto T_\sigma(w)$.

(i) Fix $\rho$ and put $j(\psi)=J_\sigma(\rho e^{i\psi})$, a $2\pi$-periodic function. Then
$\mathcal M_\sigma(\rho e^{i\psi})=\frac1{2\pi}\int_0^{2\pi}(\alpha-\pi)j(\psi+\alpha)\,d\alpha$, and integration by
parts gives
\[
 \partial_\psi\mathcal M_\sigma=\frac1{2\pi}\Bigl[(\alpha-\pi)j(\psi+\alpha)\Bigr]_0^{2\pi}
 -\frac1{2\pi}\int_0^{2\pi}j(\psi+\alpha)\,d\alpha=j(\psi)-\bar J_\sigma(\rho^2).
\]

(ii) $T_\sigma$ preserves $|w|$. With $\psi^\sharp=\psi+\mathcal M_\sigma/\bar J_\sigma$, (i) gives
$(\psi^\sharp)'=1+(J_\sigma-\bar J_\sigma)/\bar J_\sigma=J_\sigma/\bar J_\sigma>0$, and $\psi^\sharp(\psi+2\pi)=\psi^\sharp(\psi)+2\pi$
because $\mathcal M_\sigma$ is a single-valued function of $w$. So $T_\sigma$ restricts to an
orientation-preserving diffeomorphism of each circle of positive radius, and $T_\sigma(0)=0$; hence it is a
bijection of $\R^2$.

(iii) Smoothness of $(\sigma,w)\mapsto T_\sigma(w)$ was shown above. At $w=0$,
$\mathcal M_\sigma(0)=J_\sigma(0)\frac1{2\pi}\int_0^{2\pi}(\alpha-\pi)d\alpha=0$, so $\tau_\sigma(0)=0$, and
differentiating $T_\sigma(w)=\mathrm{Rot}_{\tau_\sigma(w)}w$ at $w=0$ gives
$DT_\sigma(0)=\mathrm{Rot}_0=\mathrm{Id}$. On $w\ne0$, in polar coordinates $du\wedge dv=\rho\,d\rho\wedge d\psi$,
and $\bar J_\sigma(|T_\sigma w|^2)=\bar J_\sigma(\rho^2)$, so
\[
 T_\sigma^*\bigl(\bar J_\sigma\,\rho\,d\rho\wedge d\psi\bigr)=\bar J_\sigma\,\rho\,d\rho\wedge((\psi^\sharp)'d\psi+\partial_\rho\psi^\sharp\,d\rho)
 =J_\sigma\,\rho\,d\rho\wedge d\psi .
\]
Thus $\det DT_\sigma=J_\sigma/\bar J_\sigma(|w|^2)$ on $w\ne0$; both sides are continuous, so this holds at $0$
as well (where both equal $1$), and the pull-back identity holds on $\R^2$. In particular $DT_\sigma$ is
everywhere invertible, and $T_\sigma$ is a diffeomorphism by the inverse function theorem. For the inverse,
let $(\sigma_0,x_0)\in\mathcal P\times\R^2$, $w_0=T_{\sigma_0}^{-1}(x_0)$, and choose $\rho_1>|w_0|$ and
$N_0$ as above. The map $(\sigma,w)\mapsto(\sigma,\tilde T(\sigma,w))$ is $C^\infty$ on
$N_0\times B(0,\rho_1)$ with derivative of determinant $\det D_w\tilde T(\sigma_0,w_0)>0$ at $(\sigma_0,w_0)$, so it
has a $C^\infty$ inverse $\mathcal L$ from an open neighbourhood $N$ of $(\sigma_0,x_0)$ onto an open
neighbourhood of $(\sigma_0,w_0)$. If $(\sigma,x)\in N$ with $\sigma\in\mathcal P$, then
$\mathcal L(\sigma,x)=(\sigma,w')$ with $T_\sigma(w')=x$, so $w'=T_\sigma^{-1}(x)$ by injectivity of $T_\sigma$.
Hence $(\sigma,x)\mapsto T_\sigma^{-1}(x)$ is smooth.

(iv) For $|w|\le\rho_0$ every point $\mathrm{Rot}_\alpha w$ lies in the disc, so $\bar J_\sigma(|w|^2)=j$ and
$\mathcal M_\sigma(w)=j\frac1{2\pi}\int_0^{2\pi}(\alpha-\pi)d\alpha=0$; hence $T_\sigma(w)=w$. Since $T_\sigma$ is a
bijection preserving every circle and fixing the closed disc pointwise, $T_\sigma^{-1}$ fixes it too.
\end{proof}

\begin{lemma}[Area normalisation]\label{lem:pln-area}
\lean{SymplecticPacking.Inner.Acal_hasDerivAt, SymplecticPacking.Inner.Acal_bijective,
 SymplecticPacking.Inner.Darea_bijective, SymplecticPacking.Inner.norm_Darea_sq,
 SymplecticPacking.Strip.Darea_pullback, SymplecticPacking.Strip.DareaInv_contDiff,
 SymplecticPacking.Strip.normalize_general, SymplecticPacking.Inner.Darea_eq_of_const,
 SymplecticPacking.Inner.Darea_param_smooth, SymplecticPacking.Inner.normalize_param_smooth}
\uses{def:pln-normalisations, lem:pln-rotation}
In setting (N), let $\sigma\in\mathcal P$ and put $\Omega=\{(\sigma,z):\sigma\in\mathcal P,\ \pi|z|^2<\mathcal A_\sigma(\infty)\}$.
\begin{enumerate}[label=(\roman*)]
\item $\mathcal A_\sigma(E)=\pi\int_0^E\bar J_\sigma(t)\,dt$; $\mathcal A_\sigma$ is a continuous strictly increasing
 bijection of $[0,\infty)$ onto $[0,\mathcal A_\sigma(\infty))$, so $E_\sigma$ is well defined.
\item $(\sigma,w)\mapsto\mathcal G_\sigma(w)$ is smooth and positive on $\mathcal P\times\R^2$, $\mathcal G_\sigma(0)=J_\sigma(0)$, and
 $\mathcal A_\sigma(|w|^2)=\pi|w|^2\mathcal G_\sigma(w)$.
\item $\Psi_\sigma(w)=\sqrt{\mathcal G_\sigma(w)}\,w$ is a diffeomorphism of $\R^2$ onto $D(\mathcal A_\sigma(\infty))$ with
 $\pi|\Psi_\sigma(w)|^2=\mathcal A_\sigma(|w|^2)$ and $\Psi_\sigma^*\omega_\C=\bar J_\sigma(|w|^2)\,du\wedge dv$; its
 inverse is $D_\sigma$. The map $(\sigma,z)\mapsto D_\sigma(z)$ is smooth on $\Omega$, including $z=0$;
 $|D_\sigma(z)|^2=E_\sigma(\pi|z|^2)$, $D_\sigma^*(\bar J_\sigma(|w|^2)\,du\wedge dv)=\omega_\C$,
 $DD_\sigma(0)=J_\sigma(0)^{-1/2}\,\mathrm{Id}$, and the quotient in the definition of $D_\sigma$ is
 \[
  \frac{\pi E_\sigma(\pi|z|^2)}{\pi|z|^2}=\frac1{\mathcal G_\sigma(D_\sigma(z))}\qquad(z\ne0),
 \]
 a smooth positive function of $(\sigma,z)\in\Omega$ with value $1/J_\sigma(0)$ at $z=0$.
\item (Composite.) Suppose $(\sigma,w)\mapsto F_\sigma(w)$ is $C^\infty$ on an open neighbourhood of
 $\mathcal P\times\R^2$ and each $F_\sigma$ is a diffeomorphism of $\R^2$ onto an open set $\mathcal V_\sigma$ with
 $F_\sigma^*(dx\wedge d\theta)=J_\sigma\,du\wedge dv$. Then $f_\sigma=F_\sigma\circ T_\sigma^{-1}\circ D_\sigma$ is a
 symplectic diffeomorphism of $D(\mathcal A_\sigma(\infty))$ onto $\mathcal V_\sigma$, $(\sigma,z)\mapsto f_\sigma(z)$ is smooth on
 $\Omega$, $f_\sigma$ maps the circle $\{\pi|z|^2=I\}$ onto $F_\sigma(\{|w|^2=E_\sigma(I)\})$, $f_\sigma(0)=F_\sigma(0)$ and
 $Df_\sigma(0)=J_\sigma(0)^{-1/2}DF_\sigma(0)$.
\item (Constant near the origin.) If $J_\sigma(w)=j$ for $|w|\le\rho_0$, then $\mathcal G_\sigma=j$ there,
 $\mathcal A_\sigma(E)=\pi jE$ for $E\le\rho_0^2$, $D_\sigma(z)=z/\sqrt j$ for $|z|\le\sqrt j\,\rho_0$, and in (iv)
 $f_\sigma(z)=F_\sigma(z/\sqrt j)$ for $|z|\le\sqrt j\,\rho_0$.
\end{enumerate}
\end{lemma}

\begin{proof}
(i) In polar coordinates,
\begin{align*}
 \mathcal A_\sigma(E)&=\int_0^{\sqrt E}\rho\int_0^{2\pi}J_\sigma\,d\psi\,d\rho\\
 &=\int_0^{\sqrt E}2\pi\rho\,\bar J_\sigma(\rho^2)\,d\rho
 =\pi\int_0^E\bar J_\sigma(t)\,dt.
\end{align*}
Since $J_\sigma>0$ is continuous, $\mathcal A_\sigma$ is continuous and strictly
increasing with $\mathcal A_\sigma(0)=0$, and it increases to $\mathcal A_\sigma(\infty)$.

(ii) For $w\ne0$ the real-linear map $\zeta\mapsto w\zeta$ has determinant $|w|^2$ and maps the unit disc
onto $\{|x|\le|w|\}$, so $\pi|w|^2\mathcal G_\sigma(w)=\int_{|x|\le|w|}J_\sigma(x)\,dx=\mathcal A_\sigma(|w|^2)$; at $w=0$,
$\mathcal G_\sigma(0)=J_\sigma(0)$. With $N_0$ and $\rho_1$ as in the proof of Lemma~\ref{lem:pln-rotation}, the
integrand $J(\sigma,w\zeta)$ is $C^\infty$ in $(\sigma,w,\zeta)\in N_0\times B(0,\rho_1)\times\bar B(0,1)$ and positive,
so $\mathcal G$ has a positive $C^\infty$ extension $\tilde{\mathcal G}$ to $N_0\times B(0,\rho_1)$.

(iii) In polar coordinates $w=\rho e^{i\psi}$, $\Psi_\sigma(w)=r(\rho)e^{i\psi}$ with
$\pi r(\rho)^2=\mathcal A_\sigma(\rho^2)$ by (ii). As $\rho$ runs over $[0,\infty)$, $\pi r^2$ increases strictly and
continuously from $0$ to $\mathcal A_\sigma(\infty)$; hence $\Psi_\sigma$ is a bijection of $\R^2$ onto
$D(\mathcal A_\sigma(\infty))$, and $E_\sigma(\pi|\Psi_\sigma(w)|^2)=|w|^2$. On $w\ne0$, by (i),
\[
 \Psi_\sigma^*\omega_\C=r\,dr\wedge d\psi=\tfrac1{2\pi}\,d(\pi r^2)\wedge d\psi
 =\tfrac1{2\pi}\,\pi\bar J_\sigma(\rho^2)\,2\rho\,d\rho\wedge d\psi=\bar J_\sigma(\rho^2)\,du\wedge dv ,
\]
so $\det D\Psi_\sigma(w)=\bar J_\sigma(|w|^2)$; by continuity this also holds at $w=0$, where
$D\Psi_\sigma(0)=\sqrt{\mathcal G_\sigma(0)}\,\mathrm{Id}$ and both sides equal $J_\sigma(0)$. Hence $\Psi_\sigma$ is a
diffeomorphism onto the open disc. Its inverse maps the circle $\{\pi|z|^2=I\}$ radially onto
$\{|w|^2=E_\sigma(I)\}$, i.e.\ it is $z\mapsto\sqrt{\pi E_\sigma(I)/I}\,z=D_\sigma(z)$ (and $0\mapsto0$). This gives
$|D_\sigma(z)|^2=E_\sigma(\pi|z|^2)$, the pull-back identity, and $DD_\sigma(0)=D\Psi_\sigma(0)^{-1}$. Joint smoothness:
let $(\sigma_0,z_0)\in\Omega$ and $w_0=D_{\sigma_0}(z_0)$; choose $\rho_1>|w_0|$, $N_0$ and $\tilde{\mathcal G}$ as in (ii). The map
$(\sigma,w)\mapsto(\sigma,\sqrt{\tilde{\mathcal G}(\sigma,w)}\,w)$ is $C^\infty$ on $N_0\times B(0,\rho_1)$ and its derivative at
$(\sigma_0,w_0)$ has determinant $\bar J_{\sigma_0}(|w_0|^2)>0$, so it has a $C^\infty$ inverse $\mathcal L$ on an open
neighbourhood $N$ of $(\sigma_0,z_0)$. For $(\sigma,z)\in N\cap\Omega$, $\mathcal L(\sigma,z)=(\sigma,w')$ with
$\Psi_\sigma(w')=z$, so $w'=D_\sigma(z)$ by injectivity of $\Psi_\sigma$. Finally, from $\Psi_\sigma(D_\sigma(z))=z$ we get
$|z|^2=\mathcal G_\sigma(D_\sigma(z))|D_\sigma(z)|^2$, i.e.\ $\pi E_\sigma(I)/I=|D_\sigma(z)|^2/|z|^2=1/\mathcal G_\sigma(D_\sigma(z))$, which is smooth
in $(\sigma,z)$ as a composite of smooth maps.

(iv) $D_\sigma$ is a diffeomorphism of $D(\mathcal A_\sigma(\infty))$ onto $\R^2$, $T_\sigma^{-1}$ a diffeomorphism of $\R^2$,
and $F_\sigma$ a diffeomorphism of $\R^2$ onto $\mathcal V_\sigma$. By Lemma~\ref{lem:pln-rotation}(iii) and (iii) above,
\[
 f_\sigma^*(dx\wedge d\theta)=D_\sigma^*(T_\sigma^{-1})^*\bigl(J_\sigma\,du\wedge dv\bigr)=D_\sigma^*\bigl(\bar J_\sigma\,du\wedge dv\bigr)=\omega_\C .
\]
Joint smoothness follows from that of the three factors. $D_\sigma$ maps $\{\pi|z|^2=I\}$ onto
$\{|w|^2=E_\sigma(I)\}$, which $T_\sigma^{-1}$ preserves. Finally $f_\sigma(0)=F_\sigma(0)$, and
$Df_\sigma(0)=DF_\sigma(0)\,DT_\sigma^{-1}(0)\,DD_\sigma(0)=J_\sigma(0)^{-1/2}DF_\sigma(0)$ by Lemma~\ref{lem:pln-rotation}(iii).

(v) For $|w|\le\rho_0$ and $|\zeta|\le1$, $|w\zeta|\le\rho_0$, so $\mathcal G_\sigma(w)=j$ and $\Psi_\sigma(w)=\sqrt j\,w$ on the closed disc
of radius $\rho_0$; inverting, $D_\sigma(z)=z/\sqrt j$ for $|z|\le\sqrt j\rho_0$, and $\mathcal A_\sigma(E)=\pi jE$ for
$E\le\rho_0^2$ by (ii). The last claim follows from Lemma~\ref{lem:pln-rotation}(iv).
\end{proof}

%----------------------------------------------------------------------
\subsection{The inner map}\label{sec:pln-inner}

The polar chart places the Morse rectangle into the disc of area $h$. For $y\in\mathcal Y$ let
\begin{equation}\label{eq:pln-chart}
 C_y(x,\theta)=-\sqrt{\frac{p^*(y)+x}\pi}\;e^{2\pi i\theta},\qquad
 (x,\theta)\in\mathcal C_y=\{p^*(y)+x>0,\ |\theta|<\tfrac12\}.
\end{equation}
It has polar radius $\sqrt{(p^*+x)/\pi}$ and polar angle $2\pi\theta+\pi\in(0,2\pi)$; so $C_y$ is injective,
$\pi|C_y(x,\theta)|^2=p^*(y)+x$, the image of $\mathcal C_y$ is $\C$ minus the closed positive real axis, and
$C_y^*(dp\wedge d\vartheta)=dx\wedge d\theta$, i.e.\ $C_y$ is symplectic. Its centre line $\theta=0$ is the negative real
axis, and $C_y(0)=b_y:=(-r_b(y),0)$ with $r_b(y)=\sqrt{p^*(y)/\pi}$.

\begin{definition}[The inner map]\label{def:pln-inner-map}
\lean{SymplecticPacking.Inner.Fmor, SymplecticPacking.Inner.finner, SymplecticPacking.Inner.chartC}
\uses{lem:pln-morse-inverse, def:pln-normalisations}
Let $(U,V)$ be an admissible potential pair over $\mathcal Y$ and $y\in\mathcal Y$. Apply
the formulas of Definition~\ref{def:pln-normalisations} to the density $J_y=X_y'Y_y'$ of Lemma~\ref{lem:pln-morse-inverse}
(parameters $\sigma=y\in\mathcal P=\mathcal Y$), and write $T_y$, $D_y$, $\mathcal A_y$, $E_y$ for its normalisations. In the slice case the claims about these maps are made on the disc $D(p^\dagger)$ of Definition~\ref{def:plc-data}
and are proved with the constant-tail extension above. Put
\[
 f_y=F_y\circ T_y^{-1}\circ D_y\colon D(\mathcal A_y(\infty))\to\mathcal R_y .
\]
Let $E_{\rm rect}(y)\in(0,\infty]$ be the infimum of $U_y(p)-U_{\min}(y)+V_y(\theta)$ over
$(I_U\times I_V)\setminus\bigl((0,h)\times(-\frac12,\frac12)\bigr)$ (and $+\infty$ if this set is empty), let
$I_{\rm rect}(y)=\mathcal A_y(E_{\rm rect}(y))$ (with $\mathcal A_y(+\infty)=\mathcal A_y(\infty)$), and define the
\emph{inner map}
\[
 q_{{\rm in},y}=C_y\circ f_y\colon D\bigl(I_{\rm rect}(y)\bigr)\to\C .
\]
We also put $c_x(y)=\sqrt{a_0(y)/b_0(y)}$ and $\mathbf L_y=\operatorname{diag}(c_x(y),c_x(y)^{-1})$.
\end{definition}

By \ref{it:pln-P1}, $U_y-U_{\min}$ is monotone on each side of $p^*$ and $V_y$ on each side of $0$, so
$E_{\rm rect}(y)$ is the least of the values $U_y(0)-U_{\min}$, $U_y(h)-U_{\min}$, $V_y(-\frac12)$, $V_y(\frac12)$ at those of
the points $0,h\in\R$ and $\pm\frac12$ that lie in $I_U$, respectively $I_V$; it is positive since
$p^*(y)\in(0,h)$. By Lemma~\ref{lem:pln-morse-inverse}(iv), $\mathcal A_y(E)$ is the area of the sublevel
$\{U_y-U_{\min}+V_y\le E\}$ and $\mathcal A_y(\infty)=|I_U|\,|I_V|$. For $E<E_{\rm rect}(y)$ this sublevel lies in
$(0,h)\times(-\frac12,\frac12)$, so $\mathcal A_y(E)\le h$; letting $E\uparrow E_{\rm rect}(y)$ gives $I_{\rm rect}(y)\le h$.

\begin{theorem}[The inner map]\label{thm:pln-inner}
\lean{SymplecticPacking.Branch2.Planar.MorsePair.finner_symplectic, SymplecticPacking.Branch2.Planar.MorsePair.finner_smooth,
 SymplecticPacking.Branch2.Planar.MorsePair.finner_bijective, SymplecticPacking.Inner.finner_level,
 SymplecticPacking.Branch2.Planar.MorsePair.finner_linear, SymplecticPacking.Inner.chartC_det,
 SymplecticPacking.Inner.chartC_moment, SymplecticPacking.Inner.chartC_injOn}
\uses{def:pln-inner-map, lem:pln-morse-inverse, lem:pln-rotation, lem:pln-area}
Let $(U,V)$ be an admissible potential pair. Put $\mathcal D_y^{\rm src}=\C$ for the equal-ball family and
$\mathcal D_y^{\rm src}=D(p^\dagger)$, with $p^\dagger$ of Definition~\ref{def:plc-data}, for the slice family.
\begin{enumerate}[label=(\alph*)]
\item For each $y$, $f_y$ is a symplectic embedding of $\mathcal D_y^{\rm src}$ into the Morse rectangle
 $\mathcal R_y$, and $(y,z)\mapsto f_y(z)$ is smooth on the corresponding source set. In the equal-ball case it is a diffeomorphism onto $\mathcal R_y=\R^2$. Moreover
 $f_y(0)=0$ and $Df_y(0)=\mathbf L_y$, a diagonal matrix of determinant $1$.
\item (Energy identity.) If $f_y(z)=(x,\theta)$, then
 \[
  U_y\bigl(p^*(y)+x\bigr)+V_y(\theta)=U_{\min}(y)+E_y(\pi|z|^2),
 \]
 so $f_y$ maps each circle $\{\pi|z|^2=I\}$ onto the level curve of $U_y(p^*+\cdot)-U_{\min}+V_y$ which encloses area $I$.
\item The inner map $q_{{\rm in},y}$ is a symplectic embedding of $D(I_{\rm rect}(y))\cap\mathcal D_y^{\rm src}$ into
 $D(h)\setminus[0,\infty)$, jointly smooth in $(y,z)$, with $q_{{\rm in},y}(0)=b_y$,
 $Dq_{{\rm in},y}(0)=\mathbf A_y:=DC_y(0)\,\mathbf L_y$, and, for $f_y(z)=(x,\theta)$,
 \[
  \pi|q_{{\rm in},y}(z)|^2=p^*(y)+x,\qquad
  U_y\bigl(\pi|q_{{\rm in},y}(z)|^2\bigr)+V_y(\theta)=U_{\min}(y)+E_y(\pi|z|^2).
 \]
 In particular $U_y(\pi|q_{{\rm in},y}(z)|^2)\le U_{\min}(y)+E_y(\pi|z|^2)$, with equality on the centre line $\theta=0$.
\item If the pair is exactly quadratic near the minimum, then $f_y(z)=\mathbf L_yz$ for
 $|z|\le\sqrt{j_0(y)}\,\rho_T(y)$.
\end{enumerate}
\end{theorem}

\begin{proof}
(a) For the whole-plane family, apply Lemma~\ref{lem:pln-area}(iv) in setting (N) with
$\mathcal P=\mathcal Y$, $\mathcal U=\mathcal Y\times\R^2$, the density $J_y$ and $F_y$
(Lemma~\ref{lem:pln-morse-inverse}(iv)), whose image is $\mathcal V_y=\mathcal R_y$.
For the slice family, use the constant-tail extension following
Definition~\ref{def:pln-normalisations}. Its area and rotation normalisations are smooth and have the
stated determinant identities by Lemmas~\ref{lem:pln-area} and~\ref{lem:pln-rotation}.
On $D(p^\dagger)$ they agree with the original normalisations, and their Morse images stay where the
extended density equals $J_y$. There $\det DF_y=J_y$ gives the same symplectic identity after composition.
Thus the assertion holds on each source domain. Further $f_y(0)=F_y(0)=0$
and $Df_y(0)=j_0^{-1/2}\operatorname{diag}(a_0,b_0)=\operatorname{diag}(\sqrt{a_0/b_0},\sqrt{b_0/a_0})=\mathbf L_y$.

(b) $(u,v)=T_y^{-1}D_y(z)$ satisfies $u^2+v^2=|D_y(z)|^2=E_y(\pi|z|^2)$, since $T_y^{-1}$ preserves $|w|$
(Lemmas~\ref{lem:pln-rotation}(ii) and~\ref{lem:pln-area}(iii)). Apply the level identity of
Lemma~\ref{lem:pln-morse-inverse}(iv).

(c) If $\pi|z|^2<I_{\rm rect}(y)$ then $E_y(\pi|z|^2)<E_{\rm rect}(y)$, so by (b) the point
$(p^*+x,\theta)$ lies in a sublevel below $E_{\rm rect}$, hence in $(0,h)\times(-\frac12,\frac12)$; thus
$(x,\theta)\in\mathcal C_y$ and $0<\pi|q_{{\rm in},y}(z)|^2=p^*+x<h$. The map $q_{{\rm in},y}$ is the composite of the
symplectic embedding $f_y$ and the symplectic embedding $C_y$, smooth in $(y,z)$ because $C_y$ depends smoothly on
$(y,x,\theta)$ through $p^*(y)$. The value and derivative at $0$ follow from (a), and the two displayed identities
from the moment formula for $C_y$ and (b). Equality on $\theta=0$ is $V_y(0)=0$.

(d) By Lemma~\ref{lem:pln-morse-inverse}(v), $J_y=j_0$ on $|w|\le\rho_T$, so by Lemma~\ref{lem:pln-area}(v)
$f_y(z)=F_y(z/\sqrt{j_0})$ for $|z|\le\sqrt{j_0}\rho_T$, and there $F_y(w)=(a_0u,b_0v)$; hence
$f_y(z)=j_0^{-1/2}\operatorname{diag}(a_0,b_0)z=\mathbf L_yz$.
\end{proof}

A computation from \eqref{eq:pln-chart} gives
$DC_y(0)=\operatorname{diag}\bigl(-1/(2\pi r_b(y)),\,-2\pi r_b(y)\bigr)$, so
\begin{equation}\label{eq:pln-A0}
 \mathbf A_y=\operatorname{diag}\Bigl(-\frac{c_x(y)}{2\pi r_b(y)},\,-\frac{2\pi r_b(y)}{c_x(y)}\Bigr),\qquad\det\mathbf A_y=1,
\end{equation}
a diagonal matrix with two negative entries. The comparison inequalities that turn the energy identity into a
rearrangement estimate depend on the particular potentials and are proved in \S\ref{sec:pll} and
\S\ref{sec:plc}: one bounds $\mathcal A_y(E)$ from below by the area of a rectangle, which bounds $E_y(I)$ from above.

%----------------------------------------------------------------------
\subsection{Density interpolation and the charts \texorpdfstring{$C_\lambda$}{C-lambda}}\label{sec:pln-charts}

The inner map is an embedding of a disc, not a map of $D(h)$ to itself, and it comes with no path to the
identity. Each of its two nonlinear factors has an obvious linearisation at the origin: the Morse chart
$F_y$ is deformed to the linear map $\operatorname{diag}(a_0,b_0)$ by interpolating $X_y,Y_y$ linearly with their
linear parts, and the chart $C_y$ is deformed to its affine part by the rescaling
$w\mapsto b_y+[C_y(\lambda w)-b_y]/\lambda$, $\lambda\to0$. The chart parameter $\lambda$ must start at $0$, so that
the path starts at an affine map.

\begin{definition}[Density family, rescaled charts and schedule]\label{def:pln-charts}
\lean{SymplecticPacking.Qhat.Xs, SymplecticPacking.Qhat.Ys, SymplecticPacking.Qhat.Js,
 SymplecticPacking.Qhat.Fs, SymplecticPacking.Qhat.fs, SymplecticPacking.Qhat.Lmap,
 SymplecticPacking.Qhat.rb, SymplecticPacking.Qhat.bpt, SymplecticPacking.Qhat.Clam,
 SymplecticPacking.Qhat.lambda0, SymplecticPacking.Qhat.lam}
\uses{lem:pln-morse-inverse, def:pln-normalisations, def:pln-inner-map, def:pre-flat, lem:pre-flat}
Let $(U,V)$ be an admissible potential pair over $\mathcal Y$.
\begin{enumerate}[label=(\alph*)]
\item For $y\in\mathcal Y$ and $s\in\R$ put
 \[
  X_{y,s}(u)=(1-s)a_0(y)u+sX_y(u),\quad Y_{y,s}(v)=(1-s)b_0(y)v+sY_y(v),\quad
  J_{y,s}(u,v)=X_{y,s}'(u)\,Y_{y,s}'(v),
 \]
 $F_{y,s}=(X_{y,s},Y_{y,s})$, and for $s\in[0,1]$ let $T_{y,s}$, $D_{y,s}$, $\mathcal A_{y,s}$, $E_{y,s}$ be the
 normalisations of $J_{y,s}$ (Definition~\ref{def:pln-normalisations} with parameters $(y,s)\in\mathcal Y\times[0,1]$),
 and $f_{y,s}=F_{y,s}\circ T_{y,s}^{-1}\circ D_{y,s}$ on $D(\mathcal A_{y,s}(\infty))$.
\item For $(y,\lambda,w)$ in the open set $\mathcal W=\{(y,\lambda,w)\in\mathcal Y\times\R\times\R^2:\ \lambda w\in\mathcal C_y\}$ put
 \[
  C_{y,\lambda}(w)=b_y+\int_0^1DC_y(t\lambda w)\,w\,dt .
 \]
\item For numbers $\lambda_0\in(0,\frac12]$ and $\delta_s\in(0,\frac14]$, the \emph{schedule} is
 \[
  \lambda(s)=\lambda_0\,S(2s)+(1-\lambda_0)\,S\Bigl(\frac{s-1+\delta_s}{\delta_s}\Bigr),\qquad s\in[0,1].
 \]
\end{enumerate}
\end{definition}

Since $\mathcal C_y$ is convex and contains $0$, $\lambda w\in\mathcal C_y$ implies $t\lambda w\in\mathcal C_y$ for $t\in[0,1]$,
so the integrand in (b) is defined and $C_{y,\lambda}$ is $C^\infty$ on $\mathcal W$. The schedule satisfies
$\lambda(0)=0$, $\lambda(1)=1$, $\lambda(s)=\lambda_0S(2s)\le\lambda_0$ for $s\le1-\delta_s$ (the \emph{early range}),
$\lambda=\lambda_0$ on $[\frac12,1-\delta_s]$, and $\lambda(s)\in[\lambda_0,1]$ for $s\ge1-\delta_s$ (the \emph{late range}),
because $\delta_s\le\frac14$; $\lambda$ is smooth.

\begin{lemma}[The rescaled charts]\label{lem:pln-chart-lambda}
\lean{SymplecticPacking.Qhat.Clam_eq, SymplecticPacking.Qhat.Clam_zero_eq,
 SymplecticPacking.Qhat.Clam_fderiv, SymplecticPacking.Qhat.Clam_det,
 SymplecticPacking.Qhat.Clam_injOn, SymplecticPacking.Qhat.Clam_smooth}
\uses{def:pln-charts}
For $(y,\lambda,w)\in\mathcal W$: if $\lambda\ne0$ then $C_{y,\lambda}(w)=b_y+[C_y(\lambda w)-b_y]/\lambda$; $C_{y,0}(w)=b_y+DC_y(0)w$;
$DC_{y,\lambda}(w)=DC_y(\lambda w)$, so $\det DC_{y,\lambda}=1$; and $C_{y,\lambda}$ is injective on
$\{w:\lambda w\in\mathcal C_y\}$.
\end{lemma}

\begin{proof}
For $\lambda\ne0$ the fundamental theorem of calculus gives
\[
 C_y(\lambda w)-C_y(0)=\int_0^1DC_y(t\lambda w)\lambda w\,dt.
\]
Divide by $\lambda$. The quotient formula has $w$-derivative $DC_y(\lambda w)$, and at $\lambda=0$ the formula is affine
with linear part $DC_y(0)$. For $\lambda\ne0$, $C_{y,\lambda}$ is $C_y\circ(\lambda\,\cdot)$ followed by an affine bijection,
and $C_y$ is injective on $\mathcal C_y$; $C_{y,0}$ is affine with $\det DC_y(0)=1$.
\end{proof}

Equivalently, with $g(\zeta)=(e^\zeta-1)/\zeta$ ($g(0)=1$) and $w=(u,v)$,
\[
 C_{y,\lambda}(u,v)=b_y-2\pi r_b(y)\,i\,v\,g(2\pi i\lambda v)-\frac{u\,e^{2\pi i\lambda v}}{\pi\bigl(r_b(y)+\sqrt{(p^*(y)+\lambda u)/\pi}\bigr)}
 \qquad\text{on }\mathcal W .
\]
For $\lambda\ne0$ this follows from $C_y(\lambda w)-b_y=r_b(1-e^{2\pi i\lambda v})+(r_b-\sqrt{(p^*+\lambda u)/\pi})\,e^{2\pi i\lambda v}$
and $r_b-\sqrt{(p^*+\lambda u)/\pi}=-\lambda u/\bigl(\pi(r_b+\sqrt{(p^*+\lambda u)/\pi})\bigr)$; for $\lambda=0$ it is
$b_y+DC_y(0)w$.

\begin{lemma}[The density family and the uniform energy bound]\label{lem:pln-energy}
\lean{SymplecticPacking.Qhat.Js_expand, SymplecticPacking.Qhat.Js_pos, SymplecticPacking.Qhat.Js_zero,
 SymplecticPacking.Qhat.fs_symplectic, SymplecticPacking.Qhat.fs_injective, SymplecticPacking.Qhat.fs_smooth,
 SymplecticPacking.Qhat.fs_zero_eq, SymplecticPacking.Qhat.fs_deriv_zero,
 SymplecticPacking.Qhat.fs_linear_near_zero, SymplecticPacking.Qhat.fs_eq_L_at_zero,
 SymplecticPacking.Qhat.fs_eq_finner}
\uses{def:pln-charts, lem:pln-morse-inverse, lem:pln-area}
Let $(U,V)$ be an admissible potential pair over $\mathcal Y$, $y\in\mathcal Y$, $s\in[0,1]$, and abbreviate
$a_0=a_0(y)$, $b_0=b_0(y)$, $j_0=a_0b_0$.
\begin{enumerate}[label=(\alph*)]
\item $X_{y,s},Y_{y,s}$ are $C^\infty$ and strictly increasing, and
 \[
  J_{y,s}=(1-s)^2j_0+s(1-s)\bigl(a_0Y_y'+b_0X_y'\bigr)+s^2J_y\ \ge\ (1-s)^2j_0+s^2J_y>0,\qquad J_{y,s}(0)=j_0 .
 \]
\item For $E\ge0$,
 \[
  \mathcal A_{y,s}(E)\ \ge\ \pi j_0(1-s)^2E+s^2\mathcal A_y(E)\ \ge\ \mathcal A^-_y(E):=\frac{\pi j_0E\,\mathcal A_y(E)}{\pi j_0E+\mathcal A_y(E)}\quad(E>0).
 \]
 In particular $\mathcal A_{y,s}(\infty)=\infty$ for $s<1$, and $\mathcal A_{y,1}(\infty)=\mathcal A_y(\infty)$.
\item (Uniform energy bound.) $\mathcal A^-_y$ is continuous and strictly increasing. If $\bar E>0$ and
 $\mathcal A_{y,s}(\bar E)\ge p^\dagger$ for every $s\in[0,1]$ (in both families this holds for the explicit $\bar E$ of
 Lemma~\ref{lem:pll-bounds}(d), respectively Definition~\ref{def:plc-data}), then $E_{y,s}(I)\le\bar E$ for
 $0\le I\le p^\dagger$ and $s\in[0,1]$.
\item $f_{y,s}$ is a symplectic embedding on the source domain $\mathcal D_y^{\rm src}$ of Theorem~\ref{thm:pln-inner}; the map $(y,s,z)\mapsto f_{y,s}(z)$ is smooth for $s\in[0,1]$ and $z\in\mathcal D_y^{\rm src}$. In the equal-ball case it is a diffeomorphism onto $\R^2$; $f_{y,s}(0)=0$ and $Df_{y,s}(0)=\mathbf L_y$ for every $s$;
 $f_{y,0}=\mathbf L_y$ and $f_{y,1}=f_y$. If $(u,v)=T_{y,s}^{-1}D_{y,s}(z)$, then $f_{y,s}(z)=(X_{y,s}(u),Y_{y,s}(v))$ and
 $u^2+v^2=E_{y,s}(\pi|z|^2)$.
\item If the pair is exactly quadratic near the minimum, then $J_{y,s}=j_0$ on $|w|\le\rho_T(y)$ and
 $f_{y,s}(z)=\mathbf L_yz$ for $|z|\le\sqrt{j_0}\,\rho_T(y)$, for every $s\in[0,1]$.
\end{enumerate}
\end{lemma}

\begin{proof}
(a) $X_{y,s}'=(1-s)a_0+sX_y'$ and $Y_{y,s}'=(1-s)b_0+sY_y'$ are positive for $s\in[0,1]$ (one of the two terms is
positive and the other nonnegative); expanding the product gives the formula, whose middle term is
nonnegative. $X_y'(0)=a_0$ and $Y_y'(0)=b_0$ give $J_{y,s}(0)=j_0$.

(b) Integrate (a) over $\{|w|^2\le E\}$: the first term gives $\pi j_0(1-s)^2E$, the middle term is nonnegative,
and the last gives $s^2\mathcal A_y(E)$. For $\alpha,\beta>0$ the minimum of $\alpha(1-s)^2+\beta s^2$ over $s\in\R$ is
$\alpha\beta/(\alpha+\beta)$. If $s<1$, the first term tends to $\infty$ with $E$. For $s=1$, $J_{y,1}=J_y$.

(c) For $E>0$, $1/\mathcal A^-_y(E)=1/(\pi j_0E)+1/\mathcal A_y(E)$, so positivity, continuity and strict increase of
$\mathcal A_y$ give those of $\mathcal A^-_y$. Since $\mathcal A_{y,s}$ is strictly increasing,
$\mathcal A_{y,s}(\bar E)\ge p^\dagger\ge I$ implies $E_{y,s}(I)\le\bar E$. The gate is Lemma~\ref{lem:pll-bounds}(d) for the
equal-ball family and is proved in Lemma~\ref{lem:plc-data} for the slice family.

(d) The density $(y,s,w)\mapsto J_{y,s}(w)$ is $C^\infty$ on $\mathcal U=\mathcal Y\times\R\times\R^2$ and positive on
$\mathcal P\times\R^2$ with $\mathcal P=\mathcal Y\times[0,1]$, and has the common positive constant germ. In the slice case apply the constant-tail extension on the common energy disc of radius $\sqrt{\bar E}$ and use normalization locality on $D(p^\dagger)$; in the equal-ball case setting (N) applies directly. $F_{y,s}$ is $C^\infty$ on $\mathcal U$, and for
$s\in[0,1]$ it is a diffeomorphism of $\R^2$ onto $X_{y,s}(\R)\times Y_{y,s}(\R)$ with Jacobian $J_{y,s}$; for $s<1$
both factors are onto $\R$ because the term $(1-s)a_0u$ (resp.\ $(1-s)b_0v$) is unbounded and $X_y,Y_y$ are
increasing. Lemma~\ref{lem:pln-area}(iv), restricted by locality to the stated source domain, gives everything except the last three assertions, with
$Df_{y,s}(0)=j_0^{-1/2}\operatorname{diag}(a_0,b_0)=\mathbf L_y$. For $s=0$ the density is the constant $j_0$, so
$\mathcal M=0$, $T_{y,0}=\mathrm{id}$, $D_{y,0}=j_0^{-1/2}\mathrm{id}$ and $f_{y,0}=\mathbf L_y$; for $s=1$, $F_{y,1}=F_y$ and
$f_{y,1}=f_y$. The last assertion is the computation in the proof of Theorem~\ref{thm:pln-inner}(b).

(e) For $|w|\le\rho_T$, $X_y'(u)=a_0$ and $Y_y'(v)=b_0$ (Lemma~\ref{lem:pln-morse-inverse}(v)), hence
$X_{y,s}'(u)=a_0$, $Y_{y,s}'(v)=b_0$ and $J_{y,s}(w)=j_0$; conclude as in Theorem~\ref{thm:pln-inner}(d).
\end{proof}

By (b) and (c), the energy needed to cover a disc of area $a<\mathcal A_y(\infty)$ is bounded uniformly in $s$, although
$J_{y,s}$ is close to neither $J_{y,0}$ nor $J_{y,1}$ for intermediate $s$. The quantitative input of the isotopy is collected in the following definition. All
entries are numbers, independent of $y$; each condition is required for every $y\in\mathcal Y$.

\begin{definition}[Isotopy data]\label{def:pln-isotopy-data}
\lean{SymplecticPacking.Qhat.Ec, SymplecticPacking.Qhat.DeltaP, SymplecticPacking.Qhat.Deltay,
 SymplecticPacking.Qhat.Dx, SymplecticPacking.Qhat.Dy, SymplecticPacking.Qhat.delta_s,
 SymplecticPacking.Qhat.pdag, SymplecticPacking.Qhat.Ebar, SymplecticPacking.Qhat.Xbar,
 SymplecticPacking.Qhat.Ybar, SymplecticPacking.Qhat.Pmin, SymplecticPacking.Qhat.Pmax,
 SymplecticPacking.Qhat.MC, SymplecticPacking.Qhat.Bstar, SymplecticPacking.Qhat.B3,
 SymplecticPacking.Qhat.cnew, SymplecticPacking.Qhat.Rq, SymplecticPacking.Qhat.Rsupp,
 SymplecticPacking.Qhat.Psupp, SymplecticPacking.Qhat.pon, SymplecticPacking.Qhat.pzero}
\uses{lem:pln-morse-inverse, lem:pln-energy, def:pln-charts}
Let $(U,V)$ be an admissible potential pair over $\mathcal Y$ with disc area $h$, and consider positive numbers
\[
 p_{\rm on},\ p_{\rm zero},\ p^\dagger,\ p^*_{\rm lo},\ p^*_{\rm hi},\ E_c,\ \Delta P,\ \Delta\theta,\ R_q,\ \delta_x,\ \delta_\theta,\
 \delta_s,\ \bar E,\ \bar X,\ \bar Y,\ \lambda_0,\ P_{\min},\ P_{\max},\ B_{\rm lin},\ B^*,\ c_{\rm sh},
\]
where the radial range bounds satisfy
\[
 P_{\min}\le\tfrac12\min(p^*_{\rm lo},\Delta P),\qquad P_{\max}\ge\max(\tfrac32p^*_{\rm hi},h).
\]
From these numbers we form the \emph{derived constants}
\begin{gather*}
 M_C=\frac1{2\sqrt{\pi P_{\min}}}+2\sqrt{\pi P_{\max}},\\
 R^\dagger=\sqrt{p^\dagger/\pi},\quad R_h=\sqrt{h/\pi},\quad r_b^+=\sqrt{p^*_{\rm hi}/\pi},\\
 R_{\rm supp}=\max\bigl(R^\dagger,\ r_b^++c_{\rm sh}B_{\rm lin},\ r_b^++c_{\rm sh}B^*,\ R_q\bigr),\quad P_{\rm supp}=\pi R_{\rm supp}^2,\quad
 p_{\rm bd}=\tfrac12(h+P_{\rm supp}).
\end{gather*}
The numbers are \emph{isotopy data} for the pair if, for every $y\in\mathcal Y$:
\begin{enumerate}[label=(D\arabic*)]
\item\label{it:pln-D1} $p_{\rm on}<p_{\rm zero}<p^\dagger$ and $p^*_{\rm lo}\le p^*(y)\le p^*_{\rm hi}<h$;
\item\label{it:pln-D2} (endpoint margins) for all $u,v$ with $u^2<E_c$, $v^2<E_c$:
 \[
  \Delta P\le p^*(y)+X_y(u)\le h-\Delta P,\qquad p^*(y)+X_y(u)\le\pi R_q^2<h,\qquad |Y_y(v)|\le\tfrac12-\Delta\theta ;
 \]
\item\label{it:pln-D3} (interpolation error) $|a_0(y)u-X_y(u)|\le\delta_x$ and $|b_0(y)v-Y_y(v)|\le \delta_\theta$ for $u^2<E_c$, $v^2<E_c$;
\item\label{it:pln-D4} (late range) $\delta_s\le\frac14$, $\delta_s\delta_x\le\frac12\Delta P$, $\delta_s\delta_\theta\le\frac12\Delta\theta$ and
 either $(1-\delta_s)^2\mathcal A_y(E_c)>p^\dagger$, or the
 uniform-energy alternative $\mathcal A_{y,s}(E_c)>p^\dagger$ for all $s\in[0,1]$;
 the slice construction uses the latter with $E_c=\bar E$;
\item\label{it:pln-D5} (uniform energy) $\mathcal A_{y,s}(\bar E)\ge p^\dagger$ for all $s\in[0,1]$, and
 $\max(a_0(y)|u|,|X_y(u)|)\le\bar X$, $\max(b_0(y)|v|,|Y_y(v)|)\le\bar Y$ for $u^2\le\bar E$, $v^2\le\bar E$;
\item\label{it:pln-D6} (early range) $\lambda_0\le\frac12$, $\lambda_0\bar X\le\frac12p^*_{\rm lo}$ and $\lambda_0\bar Y\le\frac14$;
\item\label{it:pln-D7} (ranges of the path)
 \begin{gather*}
  R^\dagger\max\bigl(1,|(\mathbf A_y)_{11}|,|(\mathbf A_y)_{22}|\bigr)\le B_{\rm lin},\qquad
  B^*\ge\min\Bigl(M_C\sqrt{\bar X^2+\bar Y^2},\ \frac{\bar X}{\sqrt{\pi p^*_{\rm lo}}}+2\pi r_b^+\bar Y\Bigr),\\
  c_{\rm sh}\le\min\Bigl(\frac14,\ \frac{R_h-r_b^+}{8\max(R^\dagger,B_{\rm lin},B^*)}\Bigr).
 \end{gather*}
\end{enumerate}
The schedule $\lambda(s)$ of Definition~\ref{def:pln-charts}(c) is formed with these $\lambda_0$ and $\delta_s$.
\end{definition}

The roles are as follows. The source disc of the inner family is $D(p^\dagger)$; the isotopy follows the inner
family on the prescribed core $D(h-4e_q)\subset D(p_{\rm on})$ and is cut off between the areas $p_{\rm on}$ and $p_{\rm zero}$; $R_q$ bounds the radius
of the endpoint image. $E_c$ is an energy level whose Morse coordinates are mapped by the endpoint
strictly inside the chart, with moment margin $\Delta P$ and angular margin $\Delta\theta$; $\delta_s$ is the length
of the late range, short enough that the interpolation error is at most half of each margin and that the
whole source disc $D(p^\dagger)$ stays below the level $E_c$; $\bar E,\bar X,\bar Y$ are crude bounds valid for all
$s$, used only in the early range, where the chart parameter is at most $\lambda_0$; $[P_{\min},P_{\max}]$ is the
range of radial coordinates met by the segments in Definition~\ref{def:pln-charts}(b) and $M_C$ bounds
$\|DC_y\|$ there; $B^*$ bounds the distance of the inner family from the centre $b_y$, $B_{\rm lin}$ bounds the linear
segment, and $c_{\rm sh}$ shrinks all of these to at most $(R_h-r_b^+)/8$; $D(P_{\rm supp})$ is the disc outside of which
nothing moves.

\begin{lemma}[Existence and order of the data]\label{lem:pln-data}
\lean{SymplecticPacking.Qhat.delta_s_pos, SymplecticPacking.Qhat.pdag_bounds,
 SymplecticPacking.Qhat.cnew_pos, SymplecticPacking.Qhat.Rsupp_lt, SymplecticPacking.Qhat.Psupp_lt,
 SymplecticPacking.Qhat.pon_bounds}
\uses{def:pln-isotopy-data, lem:pln-energy}
\begin{enumerate}[label=(\alph*)]
\item The two potential families have the explicit isotopy data given in
\S\ref{sec:pll-data} and Definition~\ref{def:plc-data}, respectively.
\item For any isotopy data: $0<\delta_s\le\frac14$, $p^\dagger<\mathcal A_y(E_c)<h$,
 $E_c\le E_{\rm rect}(y)$ and $p^\dagger<I_{\rm rect}(y)$ for every $y$, and
 \[
  R_{\rm supp}<R_h,\qquad\text{i.e.}\qquad p_{\rm on}<p^\dagger\le P_{\rm supp}<p_{\rm bd}<h .
 \]
\end{enumerate}
\end{lemma}

\begin{proof}
(a) In the equal-ball family the positivity and source order are
Lemma~\ref{lem:pll-source-order}, the pointwise and area gates are
Lemma~\ref{lem:pll-bounds}, and the range and shrinking constants are
Definition~\ref{def:pll-path-constants}, assembled in Lemma~\ref{lem:pll-instance}. For the slice family the
verification is Lemma~\ref{lem:plc-data}.

(b) $\delta_s>0$ is part of the data. $\mathcal A_y(E_c)>p^\dagger$ follows from either alternative of \ref{it:pln-D4}, taking $s=1$ in the uniform-energy case. If $|w|^2\le E_c'$ for
some $E_c'<E_c$, then $u^2,v^2<E_c$, so by \ref{it:pln-D2} $F_y(w)$ lies in
$[\Delta P-p^*,h-\Delta P-p^*]\times[-\frac12+\Delta\theta,\frac12-\Delta\theta]$; letting $E_c'\uparrow E_c$,
$\mathcal A_y(E_c)\le(h-2\Delta P)(1-2\Delta\theta)<h$. The same inclusion shows that every point of
$\{U_y-U_{\min}+V_y<E_c\}$ lies in $(0,h)\times(-\frac12,\frac12)$, so $E_c\le E_{\rm rect}(y)$, and then
$I_{\rm rect}(y)\ge\mathcal A_y(E_c)>p^\dagger$. Consequently $R^\dagger<R_h$. By \ref{it:pln-D1}, $r_b^+<R_h$, and
\ref{it:pln-D7} gives $r_b^++c_{\rm sh}\max(B_{\rm lin},B^*)\le r_b^++\frac18(R_h-r_b^+)<R_h$. By \ref{it:pln-D2}, $\pi R_q^2<h$.
Hence $R_{\rm supp}<R_h$, $P_{\rm supp}<h$, and $P_{\rm supp}\ge\pi(R^\dagger)^2=p^\dagger$.
\end{proof}

\begin{proposition}[The common inner family]\label{prop:pln-inner-family}
\lean{SymplecticPacking.Qhat.qinner, SymplecticPacking.Qhat.qinner_dom,
 SymplecticPacking.Qhat.qinner_smooth, SymplecticPacking.Qhat.qinner_symplectic,
 SymplecticPacking.Qhat.qinner_injOn, SymplecticPacking.Qhat.qinner_zero,
 SymplecticPacking.Qhat.qinner_one, SymplecticPacking.Qhat.qinner_zero,
 SymplecticPacking.Qhat.qinner_deriv_zero, SymplecticPacking.Qhat.A0_det,
 SymplecticPacking.Qhat.qinner_radial, SymplecticPacking.Qhat.Es_le_Ebar,
 SymplecticPacking.Qhat.Xs_le_Xbar, SymplecticPacking.Qhat.Ys_le_Ybar,
 SymplecticPacking.Qhat.late_margin, SymplecticPacking.Qhat.qinner_sub_b_le,
 SymplecticPacking.Qhat.qinner_one_norm_le}
\uses{lem:pln-energy, lem:pln-chart-lambda, def:pln-isotopy-data, lem:pln-data, thm:pln-inner}
Let $(U,V)$ be an admissible potential pair over $\mathcal Y$ with isotopy data. For $y\in\mathcal Y$, $s\in[0,1]$ and
$z\in\bar D(p^\dagger)$ write $(u,v)=T_{y,s}^{-1}D_{y,s}(z)$.
\begin{enumerate}[label=(\alph*)]
\item (A priori bounds.) $u^2+v^2=E_{y,s}(\pi|z|^2)\le\bar E$, $|X_{y,s}(u)|\le\bar X$ and $|Y_{y,s}(v)|\le\bar Y$.
\item (Late margins.) If $s\ge1-\delta_s$, then $u^2+v^2<E_c$ and
 \[
  \tfrac12\Delta P\le p^*(y)+X_{y,s}(u)\le h-\tfrac12\Delta P,\qquad |Y_{y,s}(v)|\le\tfrac12-\tfrac12\Delta\theta .
 \]
\item (The family.) For $t\in[0,1]$ the point $t\lambda(s)f_{y,s}(z)$ lies in $\mathcal C_y$ and has radial coordinate
 $p^*(y)+t\lambda(s)X_{y,s}(u)\in[P_{\min},P_{\max}]$. Hence
 \[
  q^{\rm in}_{y,s}=C_{y,\lambda(s)}\circ f_{y,s}\colon D(p^\dagger)\to\C
 \]
 is defined, $(y,s,z)\mapsto q^{\rm in}_{y,s}(z)$ is smooth on $\mathcal Y\times[0,1]\times D(p^\dagger)$, and each
 $q^{\rm in}_{y,s}$ is a symplectic embedding with
 \[
  q^{\rm in}_{y,0}(z)=b_y+\mathbf A_yz,\qquad q^{\rm in}_{y,1}=q_{{\rm in},y}|_{D(p^\dagger)},\qquad
  q^{\rm in}_{y,s}(0)=b_y,\qquad Dq^{\rm in}_{y,s}(0)=\mathbf A_y .
 \]
\item (Ranges.) $|q^{\rm in}_{y,s}(z)-b_y|\le B^*$ and $|q^{\rm in}_{y,1}(z)|\le R_q$ for $z\in D(p^\dagger)$.
\end{enumerate}
\end{proposition}

\begin{proof}
(a) By Lemma~\ref{lem:pln-energy}(d), $u^2+v^2=E_{y,s}(\pi|z|^2)$, and $E_{y,s}(\pi|z|^2)\le\bar E$ by \ref{it:pln-D5}
and monotonicity of $\mathcal A_{y,s}$. $X_{y,s}(u)$ is a convex combination of $a_0u$ and $X_y(u)$ with $u^2\le\bar E$, so
$|X_{y,s}(u)|\le\bar X$ by \ref{it:pln-D5}; likewise for $Y_{y,s}$.

(b) By Lemma~\ref{lem:pln-energy}(b) and the first alternative of \ref{it:pln-D4}, or directly by its uniform-energy alternative,
$\mathcal A_{y,s}(E_c)>p^\dagger\ge\pi|z|^2$, so $E_{y,s}(\pi|z|^2)<E_c$ and
$u^2,v^2<E_c$. By \ref{it:pln-D2}, $\Delta P\le p^*+X_y(u)\le h-\Delta P$ and $|Y_y(v)|\le\frac12-\Delta\theta$. By
\ref{it:pln-D3} and \ref{it:pln-D4}, $|X_{y,s}(u)-X_y(u)|=(1-s)|a_0u-X_y(u)|\le\delta_s\delta_x\le\frac12\Delta P$ and
$|Y_{y,s}(v)-Y_y(v)|\le\delta_s\delta_\theta\le\frac12\Delta\theta$.

(c) \emph{Early range} $s\le1-\delta_s$: $\lambda(s)\le\lambda_0$, so by (a) and \ref{it:pln-D6},
$|t\lambda X_{y,s}(u)|\le\lambda_0\bar X\le\frac12p^*_{\rm lo}\le\frac12p^*(y)$ and $|t\lambda Y_{y,s}(v)|\le\lambda_0\bar Y\le\frac14$; the radial
coordinate lies in $[\frac12p^*,\frac32p^*]\subset[P_{\min},P_{\max}]$. \emph{Late range} $s\ge1-\delta_s$: $\lambda\in[0,1]$, so
$p^*+t\lambda X_{y,s}(u)$ lies between $p^*$ and $p^*+X_{y,s}(u)\in[\frac12\Delta P,h-\frac12\Delta P]$, hence in
$[\min(p^*,\frac12\Delta P),\max(p^*,h)]\subset[P_{\min},P_{\max}]$, and $|t\lambda Y_{y,s}(v)|\le\frac12-\frac12\Delta\theta$ by (b).
In both cases the point lies in $\mathcal C_y$, so $(y,\lambda(s),f_{y,s}(z))\in\mathcal W$.

Smoothness follows from Lemma~\ref{lem:pln-energy}(d), the smoothness of $C$ on the open set $\mathcal W$
(Definition~\ref{def:pln-charts}) and of $\lambda$. Symplecticity: $\det Dq^{\rm in}_{y,s}=\det DC_y(\lambda f_{y,s})\det Df_{y,s}=1$
(Lemma~\ref{lem:pln-chart-lambda}). Injectivity: $f_{y,s}$ is injective, and $C_{y,\lambda}$ is injective on
$\{w:\lambda w\in\mathcal C_y\}$. At $s=0$, $\lambda=0$ and $f_{y,0}=\mathbf L_y$, so $q^{\rm in}_{y,0}(z)=b_y+DC_y(0)\mathbf L_yz$. At
$s=1$, $\lambda=1$, $C_{y,1}=C_y$ and $f_{y,1}=f_y$; since $p^\dagger<I_{\rm rect}(y)$ (Lemma~\ref{lem:pln-data}(b)),
$q^{\rm in}_{y,1}=q_{{\rm in},y}$ on $D(p^\dagger)$. Since $f_{y,s}(0)=0$, $C_{y,\lambda}(0)=b_y$ and $Df_{y,s}(0)=\mathbf L_y$,
$q^{\rm in}_{y,s}(0)=b_y$ and $Dq^{\rm in}_{y,s}(0)=DC_y(0)\mathbf L_y=\mathbf A_y$.

(d) Let $w=f_{y,s}(z)=(x,\theta)$ and $\lambda=\lambda(s)$; by (a), $|x|\le\bar X$, $|\theta|\le\bar Y$. \emph{First bound.} At a
point with radial coordinate $P$ the columns of $DC_y$ are orthogonal of lengths $1/(2\sqrt{\pi P})$ and $2\sqrt{\pi P}$,
so $\|DC_y\|\le M_C$ for $P\in[P_{\min},P_{\max}]$; by (c) this applies along the segment in
Definition~\ref{def:pln-charts}(b), and $|C_{y,\lambda}(w)-b_y|\le M_C|w|\le M_C\sqrt{\bar X^2+\bar Y^2}$. \emph{Second bound.} For
$\lambda>0$ put $P=p^*+\lambda x>0$ and $\varphi=2\pi\lambda\theta$; then, with $r_b=r_b(y)$,
\[
 C_y(\lambda w)-b_y=r_b\bigl(1-e^{i\varphi}\bigr)+\bigl(r_b-\sqrt{P/\pi}\bigr)e^{i\varphi},\qquad
 \bigl|r_b-\sqrt{P/\pi}\bigr|=\frac{|p^*-P|}{\pi(r_b+\sqrt{P/\pi})}\le\frac{\lambda|x|}{\sqrt{\pi p^*}},
\]
and $|1-e^{i\varphi}|\le|\varphi|$. Dividing by $\lambda$, $|C_{y,\lambda}(w)-b_y|\le|x|/\sqrt{\pi p^*}+2\pi r_b|\theta|\le
\bar X/\sqrt{\pi p^*_{\rm lo}}+2\pi r_b^+\bar Y$; for $\lambda=0$, $|DC_y(0)w|\le|x|/(2\pi r_b)+2\pi r_b|\theta|$ obeys the same
bound, since $2\pi r_b=2\sqrt{\pi p^*}\ge\sqrt{\pi p^*}$. Hence $|q^{\rm in}_{y,s}(z)-b_y|\le B^*$ by \ref{it:pln-D7}. Finally
$\pi|q^{\rm in}_{y,1}(z)|^2=p^*+X_y(u)\le\pi R_q^2$ by \ref{it:pln-D2}, because $u^2<E_c$ by (b).
\end{proof}

%----------------------------------------------------------------------
\subsection{The five-segment ordinary path}\label{sec:pln-five}

The images of the inner family stay within $B^*$ of $b_y$, but $B^*$ may exceed the distance $R_h-r_b$ from $b_y$ to
$\partial D(h)$; so $q^{\rm in}_{y,s}$ itself cannot be the isotopy. We give up area preservation for a moment: shrink by
$c_{\rm sh}$, run the inner family inside the shrunken picture, and expand back. The ambient flow and the area correction
of the next subsections apply to any path with the following properties.

\begin{definition}[Ordinary path]\label{def:pln-ordinary-path}
\lean{SymplecticPacking.Qhat.PathData, SymplecticPacking.Qhat.OrdinaryPath}
Let $\mathcal Y\subset\R^d$ be open and let $p^\dagger>0$, $0<P_{\rm supp}<h$ and $0<c_{\rm sh}\le1$. An
\emph{ordinary path} with these constants is a smooth map
$\mathfrak e\colon\mathcal Y\times\R\times D(p^\dagger)\to\C$, $(y,t,z)\mapsto\mathfrak e_{y,t}(z)$, together with a smooth function
$\mu\colon\R\to(0,\infty)$, such that for all $y\in\mathcal Y$ and $t\in\R$:
\begin{enumerate}[label=(O\arabic*)]
\item\label{it:pln-O1} $\mathfrak e_{y,t}=\mathrm{id}$ for $t\le0$ and $\mathfrak e_{y,t}=\mathfrak e_{y,1}$ for $t\ge1$;
\item\label{it:pln-O2} $\mathfrak e_{y,t}$ is injective, $\det D\mathfrak e_{y,t}(z)=\mu_t$ for all $z$, $\mu_t\ge c_{\rm sh}^2$ and $\mu_1=1$;
\item\label{it:pln-O3} $\pi|\mathfrak e_{y,t}(z)|^2\le P_{\rm supp}$ for all $z\in D(p^\dagger)$.
\end{enumerate}
\end{definition}

By \ref{it:pln-O1}, $\mu_t=1$ for $t\le0$ and $\mu_t=1$ for $t\ge1$, and $\partial_t\mathfrak e_{y,t}=0$ for $t\notin(0,1)$.

\begin{construction}[The five-segment path]\label{cons:pln-five-segment}
\lean{SymplecticPacking.Qhat.Msig, SymplecticPacking.Qhat.seg, SymplecticPacking.Qhat.clock,
 SymplecticPacking.Qhat.ordPath, SymplecticPacking.Qhat.ordPath_eq_seg, SymplecticPacking.Qhat.gammaDet}
\uses{prop:pln-inner-family, def:pln-isotopy-data, def:pre-flat, lem:pre-flat}
Let $(U,V)$ be an admissible potential pair over $\mathcal Y$ with isotopy data, and $c_{\rm sh}$ the constant of the data. For
$y\in\mathcal Y$ and $\sigma\in[0,1]$ let
$\mathbf M_{y,\sigma}=\mathrm{Rot}_{\pi\sigma}\operatorname{diag}\bigl(|(\mathbf A_y)_{11}|^\sigma,|(\mathbf A_y)_{22}|^\sigma\bigr)$, and for
$z\in D(p^\dagger)$ define the five segments
\begin{align*}
 \mathfrak e^0_{y,\sigma}(z)&=c_{\rm sh}^\sigma z, &
 \mathfrak e^1_{y,\sigma}(z)&=\sigma b_y+c_{\rm sh}z, &
 \mathfrak e^2_{y,\sigma}(z)&=b_y+c_{\rm sh}\,\mathbf M_{y,\sigma}z,\\
 \mathfrak e^3_{y,\sigma}(z)&=b_y+c_{\rm sh}\bigl(q^{\rm in}_{y,\sigma}(z)-b_y\bigr), &
 \mathfrak e^4_{y,\sigma}(z)&=b_y+c_{\rm sh}^{1-\sigma}\bigl(q^{\rm in}_{y,1}(z)-b_y\bigr).
\end{align*}
Put $\mathfrak e_{y,t}=\mathfrak e^j_{y,S(5t-j)}$ for $t\in[\frac j5,\frac{j+1}5]$, $j=0,\dots,4$, $\mathfrak e_{y,t}=\mathrm{id}$ for $t\le0$
and $\mathfrak e_{y,t}=\mathfrak e_{y,1}$ for $t\ge1$, and $\mu_t=c_{\rm sh}^{2S(5t)-2S(5t-4)}$.
\end{construction}

By \eqref{eq:pln-A0}, $\mathbf A_y=\mathrm{Rot}_\pi\operatorname{diag}(\alpha_1,\alpha_2)$ with
$\alpha_i=|(\mathbf A_y)_{ii}|$ and $\alpha_1\alpha_2=1$. Hence $\mathbf M_{y,0}=\mathrm{Id}$, $\mathbf M_{y,1}=\mathbf A_y$,
$\det\mathbf M_{y,\sigma}=(\alpha_1\alpha_2)^\sigma=1$, and $\|\mathbf M_{y,\sigma}\|\le\max(1,\alpha_1,\alpha_2)$.

\begin{lemma}[The five-segment path is an ordinary path]\label{lem:pln-five-segment}
\lean{SymplecticPacking.Qhat.ordPath_eq_seg, SymplecticPacking.Qhat.ordPath_smooth,
 SymplecticPacking.Qhat.ordPath_of_nonpos, SymplecticPacking.Qhat.ordPath_of_one_le,
 SymplecticPacking.Qhat.ordPath_injOn, SymplecticPacking.Qhat.ordPath_det,
 SymplecticPacking.Qhat.gammaDet_bounds, SymplecticPacking.Qhat.ordPath_image,
 SymplecticPacking.Qhat.Msig_zero, SymplecticPacking.Qhat.Msig_one, SymplecticPacking.Qhat.Msig_det,
 SymplecticPacking.Qhat.Msig_norm_le}
\uses{cons:pln-five-segment, prop:pln-inner-family, lem:pln-data, lem:pre-concat}
The map $\mathfrak e$ of Construction~\ref{cons:pln-five-segment}, with $\mu$, is an ordinary path with constants
$p^\dagger$, $P_{\rm supp}$, $h$, $c_{\rm sh}$, and $\mathfrak e_{y,1}=q^{\rm in}_{y,1}=q_{{\rm in},y}$ on $D(p^\dagger)$.
\end{lemma}

\begin{proof}
\emph{Matching.} $\mathfrak e^0_{y,0}=\mathrm{id}$, $\mathfrak e^0_{y,1}=c_{\rm sh}z=\mathfrak e^1_{y,0}$, $\mathfrak e^1_{y,1}=b_y+c_{\rm sh}z=\mathfrak e^2_{y,0}$,
$\mathfrak e^2_{y,1}=b_y+c_{\rm sh}\mathbf A_yz=b_y+c_{\rm sh}(q^{\rm in}_{y,0}(z)-b_y)=\mathfrak e^3_{y,0}$ (here $\lambda(0)=0$ is used),
$\mathfrak e^3_{y,1}=\mathfrak e^4_{y,0}$, and $\mathfrak e^4_{y,1}=q^{\rm in}_{y,1}$.

\emph{Smoothness.} Each segment is smooth on $\mathcal Y\times[0,1]\times D(p^\dagger)$ (Proposition~\ref{prop:pln-inner-family}(c)
for $j=3,4$). On each interval $[\frac j5,\frac{j+1}5]$ the clock $S(5t-j)$ of Lemma~\ref{lem:pre-concat} takes values in $[0,1]$, and all its
derivatives vanish at the endpoints. Hence all derivatives of $\mathfrak e$ containing at least one $t$-derivative vanish
at $t\in\{0,\frac15,\dots,1\}$ from both sides, while the purely spatial derivatives agree there because consecutive
segments match as maps. So the concatenation is smooth on $\mathcal Y\times\R\times D(p^\dagger)$
(Lemma~\ref{lem:pre-concat}), and \ref{it:pln-O1} holds.

\emph{Determinants and injectivity.} The spatial Jacobian determinants of the segments are $c_{\rm sh}^{2\sigma}$, $c_{\rm sh}^2$,
$c_{\rm sh}^2\det\mathbf M_{y,\sigma}=c_{\rm sh}^2$, $c_{\rm sh}^2\det Dq^{\rm in}_{y,\sigma}=c_{\rm sh}^2$ and $c_{\rm sh}^{2(1-\sigma)}$; with $\sigma=S(5t-j)$ they equal
$\mu_t$, which lies in $[c_{\rm sh}^2,1]$, and $\mu_1=1$. Segments $0,1,2$ are affine bijections; segments $3,4$ are the
embeddings $q^{\rm in}_{y,\sigma}$ followed by affine bijections. This is \ref{it:pln-O2}.

\emph{Ranges.} Let $|z|<R^\dagger$. Then $|\mathfrak e^0_{y,\sigma}(z)|<R^\dagger$, $|\mathfrak e^1_{y,\sigma}(z)|\le r_b^++c_{\rm sh}R^\dagger\le r_b^++c_{\rm sh}B_{\rm lin}$,
$|\mathfrak e^2_{y,\sigma}(z)|\le r_b^++c_{\rm sh}B_{\rm lin}$ by \ref{it:pln-D7} and the norm bound for $\mathbf M_{y,\sigma}$, and
$|\mathfrak e^3_{y,\sigma}(z)|\le r_b^++c_{\rm sh}B^*$ by Proposition~\ref{prop:pln-inner-family}(d). Finally
$\mathfrak e^4_{y,\sigma}(z)=(1-\beta)b_y+\beta q^{\rm in}_{y,1}(z)$ with $\beta=c_{\rm sh}^{1-\sigma}\in[c_{\rm sh},1]$ is a convex combination of $b_y$ and
$q^{\rm in}_{y,1}(z)$, so its modulus is at most $\max(r_b^+,R_q)$ (Proposition~\ref{prop:pln-inner-family}(d)). All these are
at most $R_{\rm supp}$, which is \ref{it:pln-O3}. The constants satisfy $P_{\rm supp}<h$ and $0<c_{\rm sh}\le\frac14$ by
Lemma~\ref{lem:pln-data}(b) and \ref{it:pln-D7}.
\end{proof}

Segments $2$ and $3$ stay within $(R_h-r_b^+)/8$ of $b_y$; segment $0$ stays in $\{|x|<R^\dagger\}$ and segment $1$
in $\{|x|\le r_b^++c_{\rm sh}R^\dagger\}$; the last segment uses only convexity of the disc. Since $\mu_t$ does not depend on
$z$ and $\mu_0=\mu_1=1$, the area correction of \S\ref{sec:pln-normal} acts on the core as the scaling
$z\mapsto\mu_t^{-1/2}z$ and is the identity there at $t=0$ and $t=1$.

%----------------------------------------------------------------------
\subsection{The cut-off ambient flow}\label{sec:pln-flow}

An ordinary path consists of embeddings of the source disc only. Push its velocity forward to the moving image,
cut it off in the \emph{source} variable between the areas $p_{\rm on}$ and $p_{\rm zero}$, and extend by zero. Cutting off in
the source keeps the field inside the image $\mathfrak e_{y,t}(D(p^\dagger))$, where $\mathfrak e_{y,t}^{-1}$ is defined. Fix
$0<p_{\rm on}<p_{\rm zero}<p^\dagger$ and put $r_{\rm on}=\sqrt{p_{\rm on}/\pi}$, $r_{\rm zero}=\sqrt{p_{\rm zero}/\pi}$. The radial
cutoff $\chi_{r_{\rm on},r_{\rm zero}}$ of Definition~\ref{def:pre-cutoff} is
$\chi_{r_{\rm on},r_{\rm zero}}(z)=1-S\bigl((\pi|z|^2-p_{\rm on})/(p_{\rm zero}-p_{\rm on})\bigr)$; by Lemma~\ref{lem:pre-cutoff}(a) it is
$C^\infty$, equals $1$ for $\pi|z|^2\le p_{\rm on}$ and $0$ for $\pi|z|^2\ge p_{\rm zero}$.

\begin{lemma}[The ambient flow]\label{lem:pln-ambient-flow}
\lean{SymplecticPacking.Qhat.chiCut, SymplecticPacking.Qhat.ordInv, SymplecticPacking.Qhat.ordVel,
 SymplecticPacking.Qhat.ambientField, SymplecticPacking.Qhat.ambientFlow,
 SymplecticPacking.Qhat.ambientField_good, SymplecticPacking.Qhat.ambientFlow_smooth,
 SymplecticPacking.Qhat.ambientFlow_bijective, SymplecticPacking.Qhat.ambientFlow_eq_ordPath,
 SymplecticPacking.Qhat.ambientFlow_det, SymplecticPacking.Qhat.ambientFlow_eq_id,
 SymplecticPacking.Qhat.ambientFlow_disk, SymplecticPacking.Qhat.ambientFlow_openDisk,
 SymplecticPacking.Qhat.ambientFlow_det_pos}
\uses{def:pln-ordinary-path, def:pre-cutoff, lem:pre-cutoff}
Let $\mathfrak e$ be an ordinary path and $0<p_{\rm on}<p_{\rm zero}<p^\dagger$. Let
$\mathcal O=\{(y,t,x):x\in\mathfrak e_{y,t}(D(p^\dagger))\}$ and define
\[
 W_{y,t}(x)=\chi_{r_{\rm on},r_{\rm zero}}\bigl(\mathfrak e_{y,t}^{-1}(x)\bigr)\,(\partial_t\mathfrak e_{y,t})\bigl(\mathfrak e_{y,t}^{-1}(x)\bigr)\quad\text{on }\mathcal O,
 \qquad W_{y,t}(x)=0\quad\text{otherwise}.
\]
Let $\phi_{y,t}$ be the flow of the time-dependent field $W_{y,t}$ with $\phi_{y,0}=\mathrm{id}$.
\begin{enumerate}[label=(\alph*)]
\item $W$ is $C^\infty$ on $\mathcal Y\times\R\times\C$ and vanishes for $t\notin(0,1)$ and for $\pi|x|^2>P_{\rm supp}$. Hence
 $\phi_{y,t}$ is defined for all $t$, each $\phi_{y,t}$ is a diffeomorphism of $\C$, and $(y,t,x)\mapsto\phi_{y,t}(x)$ is $C^\infty$.
\item $\phi_{y,t}(z)=\mathfrak e_{y,t}(z)$ for $\pi|z|^2\le p_{\rm on}$, and $\det D\phi_{y,t}(z)=\mu_t$ there.
\item $\phi_{y,t}=\mathrm{id}$ on $\{\pi|x|^2>P_{\rm supp}\}$; $\phi_{y,t}$ maps $\bar D(P)$ onto itself for every
 $P\ge P_{\rm supp}$ and $D(P)$ onto itself for every $P>P_{\rm supp}$, in particular $D(h)$ onto $D(h)$; and
 $\det D\phi_{y,t}>0$ everywhere. For $t\le0$, $\phi_{y,t}=\mathrm{id}$; for $t\ge1$, $\phi_{y,t}=\phi_{y,1}$.
\end{enumerate}
\end{lemma}

\begin{proof}
(a) The map $\hat{\mathfrak e}(y,t,z)=(y,t,\mathfrak e_{y,t}(z))$ is smooth and injective on the open set
$\mathcal Y\times\R\times D(p^\dagger)$, with Jacobian determinant $\mu_t>0$; by the inverse function theorem it is a
diffeomorphism onto the open set $\mathcal O$. Hence $(y,t,x)\mapsto\mathfrak e_{y,t}^{-1}(x)$ is smooth on $\mathcal O$, and so is $W$.
$W$ vanishes on $\mathcal O$ outside $K=\hat{\mathfrak e}(\mathcal Y\times[0,1]\times\bar D(p_{\rm zero}))$, because $\chi_{r_{\rm on},r_{\rm zero}}(z)=0$ for
$\pi|z|^2\ge p_{\rm zero}$ and $\partial_t\mathfrak e=0$ for $t\notin(0,1)$. The set $K$ is closed in $\mathcal Y\times\R\times\C$: if
$(y_n,t_n,\mathfrak e_{y_n,t_n}(z_n))\to(y,t,x)$ with $y\in\mathcal Y$, $t_n\in[0,1]$, $z_n\in\bar D(p_{\rm zero})$, pass to a subsequence
with $z_n\to z\in\bar D(p_{\rm zero})$; continuity gives $x=\mathfrak e_{y,t}(z)$. And $K\subset\mathcal O$ because $p_{\rm zero}<p^\dagger$.
The open sets $\mathcal O$ and the complement of $K$ cover $\mathcal Y\times\R\times\C$, and $W$ is smooth on both, so $W$ is
smooth. By \ref{it:pln-O3}, $W_{y,t}(x)=0$ for $\pi|x|^2>P_{\rm supp}$. A smooth time-dependent vector field depending smoothly
on parameters and vanishing outside a fixed bounded set of $x$ has a global flow by diffeomorphisms, smooth in
$(y,t,x)$.

(b) For $\pi|z|^2\le p_{\rm on}$, the curve $t\mapsto\mathfrak e_{y,t}(z)$ satisfies
$\frac d{dt}\mathfrak e_{y,t}(z)=\chi_{r_{\rm on},r_{\rm zero}}(z)\,\partial_t\mathfrak e_{y,t}(z)=W_{y,t}(\mathfrak e_{y,t}(z))$ and $\mathfrak e_{y,0}(z)=z$; by uniqueness of
solutions it is $t\mapsto\phi_{y,t}(z)$. On the open disc $D(p_{\rm on})$ the two maps agree, so their derivatives agree and
$\det D\phi_{y,t}=\mu_t$; on its boundary circle this follows by continuity of $\det D\phi_{y,t}$.

(c) Where $W_{y,t}$ vanishes for all $t$ the constant curve is a solution, so $\phi_{y,t}=\mathrm{id}$ on $\{\pi|x|^2>P_{\rm supp}\}$.
Since $\phi_{y,t}$ is a bijection of $\C$ fixing every point outside $\bar D(P_{\rm supp})$, it maps $\bar D(P_{\rm supp})$ onto itself,
hence $\bar D(P)$ and $D(P)$ onto themselves for $P\ge P_{\rm supp}$, respectively $P>P_{\rm supp}$; $h>P_{\rm supp}$. The determinant
$\det D\phi_{y,t}$ never vanishes, is continuous in $t$ and equals $1$ at $t=0$. Since $W_{y,t}=0$ for $t\notin(0,1)$, the
flow is constant in $t$ on $(-\infty,0]$ and on $[1,\infty)$.
\end{proof}

%----------------------------------------------------------------------
\subsection{Area normalisation of the ambient flow}\label{sec:pln-normal}

The flow $\phi_{y,t}$ has the right support and the right endpoint on the core, but it does not preserve area. We
correct it with the normalisations of \S\ref{sec:pln-normalisation-maps}, applied to
$G_{y,t}(w)=\phi_{y,t}(w/\sqrt{\mu_t})$, whose density is $1$ near the origin for every $t$; this lets
Lemma~\ref{lem:pln-area}(v) apply with one radius for the whole family.

\begin{lemma}[Normalisation of the whole path]\label{lem:pln-normalisation}
\lean{SymplecticPacking.Qhat.Jphi, SymplecticPacking.Qhat.Gphi, SymplecticPacking.Qhat.Jhat,
 SymplecticPacking.Inner.Acal, SymplecticPacking.Qhat.normalize, SymplecticPacking.Qhat.qhat,
 SymplecticPacking.Qhat.Jhat_eq, SymplecticPacking.Qhat.Jhat_pos, SymplecticPacking.Qhat.Jhat_smooth,
 SymplecticPacking.Qhat.Acal_eq_of_image, SymplecticPacking.Inner.Acal_bijective, SymplecticPacking.Qhat.Jhat_core,
 SymplecticPacking.Inner.Acal_eq_of_le, SymplecticPacking.Qhat.Jhat_collar, SymplecticPacking.Qhat.Acal_collar,
 SymplecticPacking.Qhat.qhat_smooth, SymplecticPacking.Qhat.qhat_symplectic,
 SymplecticPacking.Qhat.qhat_bijective, SymplecticPacking.Qhat.qhat_disk,
 SymplecticPacking.Qhat.qhat_openDisk, SymplecticPacking.Qhat.qhat_core,
 SymplecticPacking.Qhat.qhat_common_core, SymplecticPacking.Qhat.qhat_collar,
 SymplecticPacking.Qhat.qhat_zero, SymplecticPacking.Qhat.qhat_of_nonpos,
 SymplecticPacking.Qhat.qhat_of_one_le}
\uses{lem:pln-ambient-flow, lem:pln-rotation, lem:pln-area}
Let $\mathfrak e$ be an ordinary path, $0<p_{\rm on}<p_{\rm zero}<p^\dagger$, and $\phi$ the ambient flow of
Lemma~\ref{lem:pln-ambient-flow}. For $(y,t)\in\mathcal Y\times\R$ put
\[
 G_{y,t}(w)=\phi_{y,t}\bigl(w/\sqrt{\mu_t}\bigr),\qquad \hat J_{y,t}(w)=\det DG_{y,t}(w)=\frac{\det D\phi_{y,t}(w/\sqrt{\mu_t})}{\mu_t},
\]
let $\hat T_{y,t}$, $\hat D_{y,t}$ be the rotation and area normalisations of $\hat J_{y,t}$ (parameters $(y,t)$), and define
\[
 \mathcal N_{y,t}=\mu_t^{-1/2}\,\hat T_{y,t}^{-1}\circ\hat D_{y,t},\qquad
 \hat q_{y,t}=\phi_{y,t}\circ\mathcal N_{y,t}=G_{y,t}\circ\hat T_{y,t}^{-1}\circ\hat D_{y,t}.
\]
\begin{enumerate}[label=(\alph*)]
\item $\hat J$ is $C^\infty$ and positive on $\mathcal Y\times\R\times\C$; $\hat J_{y,t}(w)=1$ for $\pi|w|^2\le\mu_tp_{\rm on}$, in particular for
 $|w|\le\rho_c=c_{\rm sh}\sqrt{p_{\rm on}/\pi}$; $\hat J_{y,t}(w)=1/\mu_t$ for $\pi|w|^2>\mu_tP_{\rm supp}$; and
 $\mathcal A_{\hat J_{y,t}}(\mu_tP/\pi)=\operatorname{area}\phi_{y,t}(\bar D(P))$, which equals $\mu_tP$ for $P\le p_{\rm on}$ and $P$ for
 $P\ge P_{\rm supp}$. In particular $\mathcal A_{\hat J_{y,t}}(\infty)=\infty$.
\item Each $\hat q_{y,t}$ is a symplectic diffeomorphism of $\C$; it maps $\bar D(P)$ onto itself for every $P\ge P_{\rm supp}$ and
 $D(P)$ onto itself for every $P>P_{\rm supp}$, in particular $D(h)$ and $\bar D(h)$.
\item (Core.) $\hat q_{y,t}(z)=\mathfrak e_{y,t}(z/\sqrt{\mu_t})$ for $\pi|z|^2\le\mu_tp_{\rm on}$, in particular on the common core
 $\pi|z|^2\le c_{\rm sh}^2p_{\rm on}$.
\item (Collar.) $\hat q_{y,t}(z)=z$ for $\pi|z|^2>P_{\rm supp}$.
\item (Smoothness.) $(y,t,z)\mapsto\hat q_{y,t}(z)$ and $(y,t,x)\mapsto\hat q_{y,t}^{-1}(x)$ are $C^\infty$ on $\mathcal Y\times\R\times\C$,
 including $z=0$; $\hat q_{y,t}=\mathrm{id}$ for $t\le0$ and $\hat q_{y,t}=\hat q_{y,1}$ for $t\ge1$.
\item (Endpoint.) $\hat q_{y,1}=\mathfrak e_{y,1}$ on $\bar D(p_{\rm on})$.
\end{enumerate}
\end{lemma}

\begin{proof}
(a) $\det D\phi_{y,t}>0$ is smooth in $(y,t,x)$ (Lemma~\ref{lem:pln-ambient-flow}), and $\mu_t\ge c_{\rm sh}^2>0$ is smooth; the
formula for $\hat J$ is the chain rule. If $\pi|w|^2\le\mu_tp_{\rm on}$, then $x=w/\sqrt{\mu_t}$ has $\pi|x|^2\le p_{\rm on}$, where
$\det D\phi_{y,t}=\mu_t$ (Lemma~\ref{lem:pln-ambient-flow}(b)); so $\hat J_{y,t}(w)=1$, and $\mu_t\ge c_{\rm sh}^2$ gives the uniform
radius $\rho_c$. If $\pi|w|^2>\mu_tP_{\rm supp}$, then $\phi_{y,t}=\mathrm{id}$ near $w/\sqrt{\mu_t}$, so $\hat J_{y,t}(w)=1/\mu_t$. Since
$G_{y,t}$ is a diffeomorphism with Jacobian $\hat J_{y,t}$, the change of variables formula gives
$\mathcal A_{\hat J_{y,t}}(\mu_tP/\pi)=\operatorname{area}G_{y,t}(\{\pi|w|^2\le\mu_tP\})=\operatorname{area}\phi_{y,t}(\bar D(P))$. For $P\le p_{\rm on}$,
$\phi_{y,t}=\mathfrak e_{y,t}$ on $\bar D(P)$ with constant Jacobian $\mu_t$, so the area is $\mu_tP$; for $P\ge P_{\rm supp}$, $\phi_{y,t}$
maps $\bar D(P)$ onto itself (Lemma~\ref{lem:pln-ambient-flow}(c)), so the area is $P$.

Setting (N) applies with $\mathcal P=\mathcal Y\times\R$ (open), $\mathcal U=\mathcal P\times\R^2$ and $J=\hat J$, which equals $1$ on
$|w|\le\rho_c$ for all parameters. By Lemmas~\ref{lem:pln-rotation} and~\ref{lem:pln-area}, $\hat T$, $\hat T^{-1}$ and $\hat D$ are
smooth in all variables, $\hat T_{y,t}=\hat T_{y,t}^{-1}=\mathrm{id}$ on $|w|\le\rho_c$, $\hat D_{y,t}$ is a diffeomorphism of $\C$, and
$\hat q_{y,t}=G_{y,t}\circ\hat T_{y,t}^{-1}\circ\hat D_{y,t}$ is a symplectic diffeomorphism of $\C$ onto $G_{y,t}(\C)=\C$
(Lemma~\ref{lem:pln-area}(iv) with $F=G_{y,t}$).

(d) Let $I=\pi|z|^2>P_{\rm supp}$. By (a), $\mathcal A_{\hat J_{y,t}}(\mu_tI/\pi)=I$, so $|\hat D_{y,t}(z)|^2=\mu_tI/\pi$ and, $\hat D$ being
radial, $\hat D_{y,t}(z)=\sqrt{\mu_t}\,z$. On the circle through $w=\sqrt{\mu_t}z$ the density $\hat J_{y,t}$ is the constant $1/\mu_t$,
so $\mathcal M_{\hat J_{y,t}}(w)=0$ and $\hat T_{y,t}^{-1}(w)=w$. Hence $\mathcal N_{y,t}(z)=z$ and $\hat q_{y,t}(z)=\phi_{y,t}(z)=z$.

(b) $\hat q_{y,t}$ is a bijection of $\C$ fixing every point outside $\bar D(P_{\rm supp})$ by (d), so it maps $\bar D(P_{\rm supp})$, hence
$\bar D(P)$ for $P\ge P_{\rm supp}$ and $D(P)$ for $P>P_{\rm supp}$, onto itself.

(c) Let $\rho=\sqrt{\mu_tp_{\rm on}/\pi}$. By (a), $\hat J_{y,t}=1$ on $|w|\le\rho$, so Lemma~\ref{lem:pln-area}(v) gives
$\hat q_{y,t}(z)=G_{y,t}(z)=\phi_{y,t}(z/\sqrt{\mu_t})$ for $|z|\le\rho$; since $\pi|z/\sqrt{\mu_t}|^2\le p_{\rm on}$, this equals
$\mathfrak e_{y,t}(z/\sqrt{\mu_t})$ by Lemma~\ref{lem:pln-ambient-flow}(b).

(e) Joint smoothness of $\hat q$ was shown in the proof of (a) (all factors are smooth in $(y,t,\cdot)$). For $t\le0$,
$\phi_{y,t}=\mathrm{id}$ and $\mu_t=1$, so $\hat J=1$, $\hat T=\hat D=\mathrm{id}$ and $\hat q_{y,t}=\mathrm{id}$; for $t\ge1$, $\phi_{y,t}=\phi_{y,1}$ and
$\mu_t=1$, so $\hat q_{y,t}=\hat q_{y,1}$. The map $(y,t,z)\mapsto(y,t,\hat q_{y,t}(z))$ is a smooth bijection of
$\mathcal Y\times\R\times\C$ with Jacobian determinant $\det D\hat q_{y,t}=1$; by the inverse function theorem its inverse,
$(y,t,x)\mapsto(y,t,\hat q_{y,t}^{-1}(x))$, is smooth.

(f) is (c) for $t=1$, where $\mu_1=1$.
\end{proof}

%----------------------------------------------------------------------
\subsection{The Hamiltonian}\label{sec:pln-ham}

A symplectic isotopy of the plane is Hamiltonian: its velocity field $\xi$ satisfies $d(\iota_\xi\omega_\C)=0$, and a closed
one-form on the plane is exact. The primitive is written down by the radial (Poincar\'e) formula in the source
variable, and the additive constant is fixed by evaluation at a point of the collar, where the isotopy is the identity
and the primitive is constant.

\begin{construction}[The Hamiltonian]\label{cons:pln-hamiltonian}
\lean{SymplecticPacking.Qhat.aq, SymplecticPacking.Qhat.hq, SymplecticPacking.Qhat.pbd,
 SymplecticPacking.Qhat.zbd, SymplecticPacking.Qhat.qhatInv, SymplecticPacking.Qhat.Hq,
 SymplecticPacking.Qhat.Qham}
\uses{lem:pln-normalisation}
For an ordinary path with normalised isotopy $\hat q$, $y\in\mathcal Y$, $t\in\R$ and $z,w\in\C$, put
\begin{gather*}
 \mathfrak a_{y,t}(z)[w]=\omega_\C\bigl(\partial_t\hat q_{y,t}(z),\,D\hat q_{y,t}(z)w\bigr),\qquad
 \mathfrak h_{y,t}(z)=\int_0^1\mathfrak a_{y,t}(rz)[z]\,dr,\\
 z_{\rm bd}=\bigl(\sqrt{p_{\rm bd}/\pi},0\bigr),\quad p_{\rm bd}=\tfrac12(h+P_{\rm supp}),\qquad
 Q_{y,t}(x)=-\bigl[\mathfrak h_{y,t}-\mathfrak h_{y,t}(z_{\rm bd})\bigr]\bigl(\hat q_{y,t}^{-1}(x)\bigr).
\end{gather*}
\end{construction}

\begin{proposition}[$Q$ generates $\hat q$]\label{prop:pln-hamiltonian}
\lean{SymplecticPacking.Qhat.aq_closed, SymplecticPacking.Qhat.hq_deriv,
 SymplecticPacking.Qhat.Qham_smooth, SymplecticPacking.Qhat.Qham_support,
 SymplecticPacking.Qhat.qhat_hamiltonian, SymplecticPacking.Qhat.qhat_interface}
\uses{cons:pln-hamiltonian, lem:pln-normalisation}
$Q$ is $C^\infty$ on $\mathcal Y\times\R\times\C$; $Q_{y,t}$ and all its derivatives vanish on $\{\pi|x|^2>P_{\rm supp}\}$ for every
$(y,t)$, and $Q_{y,t}=0$ for $t\notin(0,1)$. With $\xi_{y,t}=(\partial_t\hat q_{y,t})\circ\hat q_{y,t}^{-1}$,
\[
 \iota_{\xi_{y,t}}\omega_\C=-dQ_{y,t},\qquad \partial_t\hat q_{y,t}=\xi_{y,t}\circ\hat q_{y,t},
\]
so $\hat q_{y,\cdot}$ is the isotopy generated by $Q_{y,\cdot}$ in the convention $\iota_X\omega_\C=-dQ$.
\end{proposition}

\begin{proof}
By definition $\mathfrak a_{y,t}=\hat q_{y,t}^*(\iota_{\xi_{y,t}}\omega_\C)$. Differentiating $\hat q_{y,t}^*\omega_\C=\omega_\C$ in $t$ gives
$\hat q_{y,t}^*\mathcal L_{\xi_{y,t}}\omega_\C=0$, so $d\iota_{\xi_{y,t}}\omega_\C=\mathcal L_{\xi_{y,t}}\omega_\C=0$ and $\mathfrak a_{y,t}$ is closed. The plane
is star-shaped about $0$, so the radial primitive satisfies $d\mathfrak h_{y,t}=\mathfrak a_{y,t}$; it is smooth in $(y,t,z)$ because the
integrand is. Then $-dQ_{y,t}=(\hat q_{y,t}^{-1})^*d\mathfrak h_{y,t}=\iota_{\xi_{y,t}}\omega_\C$. On the connected set
$\{\pi|z|^2>P_{\rm supp}\}$, $\hat q_{y,t}=\mathrm{id}$ for all $t$ (Lemma~\ref{lem:pln-normalisation}(d)), so $\partial_t\hat q_{y,t}=0$ and
$\mathfrak a_{y,t}=0$ there; hence $\mathfrak h_{y,t}$ is constant on this set, equal to $\mathfrak h_{y,t}(z_{\rm bd})$ because
$P_{\rm supp}<p_{\rm bd}$. Therefore $Q_{y,t}$ vanishes identically on $\hat q_{y,t}(\{\pi|z|^2>P_{\rm supp}\})=\{\pi|x|^2>P_{\rm supp}\}$, with all
derivatives. For $t\notin(0,1)$, $\partial_t\hat q_{y,t}=0$, so $\mathfrak a_{y,t}=0$ and $Q_{y,t}=0$. Smoothness of $Q$ follows from
that of $\mathfrak h$ and of $(y,t,x)\mapsto\hat q_{y,t}^{-1}(x)$ (Lemma~\ref{lem:pln-normalisation}(e)).
\end{proof}

The radial primitive gives $\hat q_{y,t}^*(\iota_{\xi_{y,t}}\omega_\C)=+d\mathfrak h_{y,t}$ in the source. The sign of $Q$ matches the
fibre convention $\iota_X\omega_\C=-dQ$ of \Msec{sec:cmp}, which is the convention used for the suspension in
\Msec{sec:hdl}.

%----------------------------------------------------------------------
\subsection{The whole-disc isotopy}\label{sec:pln-isotopy}

\begin{theorem}[Whole-disc Hamiltonian isotopy with prescribed endpoint on the core]\label{thm:pln-isotopy}
\lean[needed]{SymplecticPacking.Qhat.qhat_interface, SymplecticPacking.Qhat.qhat_one_eq,
 SymplecticPacking.Qhat.qhat_disk, SymplecticPacking.Qhat.qhat_openDisk,
 SymplecticPacking.Qhat.qhat_of_nonpos, SymplecticPacking.Qhat.qhat_of_one_le, SymplecticPacking.Qhat.qhat_core,
 SymplecticPacking.Qhat.qhatP, SymplecticPacking.Qhat.qhatP_spec}
\uses{prop:pln-inner-family, lem:pln-five-segment, lem:pln-ambient-flow, lem:pln-normalisation, prop:pln-hamiltonian, thm:pln-inner}
Let $(U,V)$ be an admissible potential pair over $\mathcal Y$ with disc area $h$, and let isotopy data be given
(Definition~\ref{def:pln-isotopy-data}). Let $\hat q$ and $Q$ be obtained from the five-segment path of
Construction~\ref{cons:pln-five-segment} by Lemma~\ref{lem:pln-normalisation}, with the levels $p_{\rm on}<p_{\rm zero}$ of the data,
and Construction~\ref{cons:pln-hamiltonian}. Then $P_{\rm supp}<h$, and for every $y\in\mathcal Y$:
\begin{enumerate}[label=(\alph*)]
\item each $\hat q_{y,t}$ is a symplectic diffeomorphism of $\C$ which is the identity on $\{\pi|z|^2>P_{\rm supp}\}$ and maps
 $\bar D(P)$ onto itself for every $P\ge P_{\rm supp}$, in particular $\bar D(h)$ and $D(h)$;
\item $\hat q_{y,t}=\mathrm{id}$ for $t\le0$ and $\hat q_{y,t}=\hat q_{y,1}$ for $t\ge1$;
\item (same endpoint on the core) $\hat q_{y,1}(z)=q_{{\rm in},y}(z)$ for $\pi|z|^2\le p_{\rm on}$;
\item $Q_{y,t}$ vanishes with all derivatives on $\{\pi|x|^2>P_{\rm supp}\}$ and for $t\notin(0,1)$, and generates $\hat q_{y,t}$:
 $\partial_t\hat q_{y,t}=\xi_{y,t}\circ\hat q_{y,t}$ with $\iota_{\xi_{y,t}}\omega_\C=-dQ_{y,t}$;
\item (core formula) $\hat q_{y,t}(z)=\mathfrak e_{y,t}(z/\sqrt{\mu_t})$ for $\pi|z|^2\le\mu_tp_{\rm on}$.
\end{enumerate}
\end{theorem}

\begin{proof}
By Lemma~\ref{lem:pln-five-segment} the five-segment path is an ordinary path with constants
$p^\dagger$, $P_{\rm supp}$, $h$, $c_{\rm sh}$; $p_{\rm on}<p_{\rm zero}<p^\dagger$ by \ref{it:pln-D1}, and $P_{\rm supp}<h$ by
Lemma~\ref{lem:pln-data}(b). Items (a), (b) and (e) are
Lemma~\ref{lem:pln-normalisation}(b)--(e); (d) is Proposition~\ref{prop:pln-hamiltonian}. For (c),
Lemma~\ref{lem:pln-normalisation}(f) gives $\hat q_{y,1}=\mathfrak e_{y,1}$ on $\bar D(p_{\rm on})$, and
$\mathfrak e_{y,1}=q^{\rm in}_{y,1}=q_{{\rm in},y}$ on $D(p^\dagger)\supset\bar D(p_{\rm on})$ by Lemma~\ref{lem:pln-five-segment}.
\end{proof}

On the annulus $\{p_{\rm on}<\pi|z|^2\le h\}$ we use only $\hat q_{y,1}(\bar D(h))=\bar D(h)$:

\begin{corollary}[Comparison on the whole closed disc]\label{thm:pln-full-disc}
\lean[needed]{SymplecticPacking.Qhat.R_q_le, SymplecticPacking.Qhat.endpoint_comparison, SymplecticPacking.Qhat.full_comparison}
\uses{thm:pln-isotopy}
In the situation of Theorem~\ref{thm:pln-isotopy}, let $y\in\mathcal Y$, let $\Upsilon\colon[0,h]\to\R$ be any function,
$g\colon[0,h]\to\R$ nondecreasing and $\eps\ge0$, and suppose that $\Upsilon(\pi|q_{{\rm in},y}(z)|^2)\le g(\pi|z|^2)+\eps$ whenever
$\pi|z|^2\le p_{\rm on}$, and that $\sup_{[0,h]}\Upsilon\le g(p_{\rm on})+\eps$. Then
\[
 \Upsilon\bigl(\pi|\hat q_{y,1}(z)|^2\bigr)\le g(\pi|z|^2)+\eps\qquad\text{for all }z\in\bar D(h).
\]
\end{corollary}

\begin{proof}
For $\pi|z|^2\le p_{\rm on}$ use Theorem~\ref{thm:pln-isotopy}(c). For $p_{\rm on}<\pi|z|^2\le h$, $\hat q_{y,1}(z)\in\bar D(h)$ by
Theorem~\ref{thm:pln-isotopy}(a), so $\Upsilon(\pi|\hat q_{y,1}(z)|^2)\le\sup\Upsilon\le g(p_{\rm on})+\eps\le g(\pi|z|^2)+\eps$.
\end{proof}

All constructions above were made with $y$ as a variable, and the data are numbers independent of $y$. The
following theorem records the consequences used for families.

\begin{theorem}[Parametrised version]\label{thm:pln-param}
\uses{thm:pln-isotopy, lem:pln-normalisation, prop:pln-hamiltonian, prop:pln-inner-family}
In the situation of Theorem~\ref{thm:pln-isotopy}:
\begin{enumerate}[label=(\alph*)]
\item the maps $(y,t,z)\mapsto\hat q_{y,t}(z)$, $(y,t,x)\mapsto\hat q_{y,t}^{-1}(x)$ and $(y,t,x)\mapsto Q_{y,t}(x)$ are $C^\infty$ on
 $\mathcal Y\times\R\times\C$, and $(y,s,z)\mapsto q^{\rm in}_{y,s}(z)$ is smooth on $\mathcal Y\times[0,1]\times D(p^\dagger)$; the
 numbers $p_{\rm on}$, $P_{\rm supp}$, $c_{\rm sh}$, and hence the collar $\{\pi|z|^2>P_{\rm supp}\}$, the core $\bar D(p_{\rm on})$ and the
 common core $\bar D(c_{\rm sh}^2p_{\rm on})$, do not depend on $y$;
\item if $\mathcal Y_0\subset\mathcal Y$ is compact, then the derivatives of order $\ge1$ of $\hat q$ and $\hat q^{-1}$, and all
 derivatives (including order $0$) of $\hat q-\mathrm{id}$, $\hat q^{-1}-\mathrm{id}$ and $Q$, are bounded on
 $\mathcal Y_0\times\R\times\C$.
\end{enumerate}
\end{theorem}

\begin{proof}
(a) is Lemma~\ref{lem:pln-normalisation}(e), Proposition~\ref{prop:pln-hamiltonian} and
Proposition~\ref{prop:pln-inner-family}(c). (b) The maps are smooth on $\mathcal Y\times\R\times\C$ and constant in $t$ outside
$[0,1]$, and $\hat q_{y,t}-\mathrm{id}$, $\hat q^{-1}_{y,t}-\mathrm{id}$ and $Q_{y,t}$ vanish outside $\bar D(P_{\rm supp})$. Hence each of
the listed derivatives is a continuous function on $\mathcal Y_0\times\R\times\C$ which, outside the compact set
$\mathcal Y_0\times[0,1]\times\bar D(P_{\rm supp})$, coincides with its value at $t=0$ or $t=1$, or equals $0$, or (for first
derivatives of $\hat q$ and $\hat q^{-1}$) the identity matrix; so it is bounded.
\end{proof}

\begin{remark}[Explicitness]\label{rem:pln-explicit}
Every object of this section is a defined term in the sense of \Mref{Definition}{def:intro-levels}: finite expressions in the
potentials and their derivatives, integrals over compact domains ($\bar J$, $\mathcal M$, $\mathcal G$, the charts
$C_{y,\lambda}$, the radial primitive), inverses of explicit bijections with nonvanishing derivative ($X_y$, $Y_y$,
$\mathcal A_{y,s}$, $\Psi_\sigma$, $T_\sigma$, $(y,t,z)\mapsto(y,t,\mathfrak e_{y,t}(z))$, $\hat q_{y,t}$), and the flow of an explicit compactly supported vector field. The base point
$z_{\rm bd}$, the schedule $\lambda(s)$, the clocks $S(5t-j)$ of Lemma~\ref{lem:pre-concat} and the time intervals
$[\frac j5,\frac{j+1}5]$ are fixed. The isotopy data are the numbers of \S\ref{sec:pll} and \S\ref{sec:plc}, given by finite expressions; square roots and compact extrema occur in both constructions.
\end{remark}

Theorem~\ref{thm:pln-isotopy}, Corollary~\ref{thm:pln-full-disc} and Theorem~\ref{thm:pln-param} supply exactly the three
properties of the fibre isotopy used in \Msec{sec:hdl}.
